% Generated reading copy of docs/PATHWAYS.md, which stays the authority.
% Do not edit by hand: regenerate it whenever PATHWAYS.md changes, with
% pandoc 3.9 (gfm+tex_math_dollars, --toc, a Lua filter that sizes table
% columns by their longest cell, and the \DeclareUnicodeCharacter / \texttt
% preamble lines below). Build: pdflatex, twice.
% Options for packages loaded elsewhere
\PassOptionsToPackage{unicode}{hyperref}
\PassOptionsToPackage{hyphens}{url}
\PassOptionsToPackage{dvipsnames,svgnames,x11names}{xcolor}
\documentclass[
  11pt,
]{article}
\usepackage{xcolor}
\usepackage[margin=2.5cm]{geometry}
\usepackage{amsmath,amssymb}
\setcounter{secnumdepth}{-\maxdimen} % remove section numbering
\usepackage{iftex}
\ifPDFTeX
  \usepackage[T1]{fontenc}
  \usepackage[utf8]{inputenc}
  \usepackage{textcomp} % provide euro and other symbols
\else % if luatex or xetex
  \usepackage{unicode-math} % this also loads fontspec
  \defaultfontfeatures{Scale=MatchLowercase}
  \defaultfontfeatures[\rmfamily]{Ligatures=TeX,Scale=1}
\fi
\usepackage{lmodern}
\ifPDFTeX\else
  % xetex/luatex font selection
\fi
% Use upquote if available, for straight quotes in verbatim environments
\IfFileExists{upquote.sty}{\usepackage{upquote}}{}
\IfFileExists{microtype.sty}{% use microtype if available
  \usepackage[]{microtype}
  \UseMicrotypeSet[protrusion]{basicmath} % disable protrusion for tt fonts
}{}
\makeatletter
\@ifundefined{KOMAClassName}{% if non-KOMA class
  \IfFileExists{parskip.sty}{%
    \usepackage{parskip}
  }{% else
    \setlength{\parindent}{0pt}
    \setlength{\parskip}{6pt plus 2pt minus 1pt}}
}{% if KOMA class
  \KOMAoptions{parskip=half}}
\makeatother
\usepackage{longtable,booktabs,array}
\newcounter{none} % for unnumbered tables
\usepackage{calc} % for calculating minipage widths
% Correct order of tables after \paragraph or \subparagraph
\usepackage{etoolbox}
\makeatletter
\patchcmd\longtable{\par}{\if@noskipsec\mbox{}\fi\par}{}{}
\makeatother
% Allow footnotes in longtable head/foot
\IfFileExists{footnotehyper.sty}{\usepackage{footnotehyper}}{\usepackage{footnote}}
\makesavenoteenv{longtable}
\setlength{\emergencystretch}{3em} % prevent overfull lines
\providecommand{\tightlist}{%
  \setlength{\itemsep}{0pt}\setlength{\parskip}{0pt}}
\DeclareUnicodeCharacter{2032}{\ensuremath{'}}
\DeclareUnicodeCharacter{2033}{\ensuremath{''}}
\DeclareUnicodeCharacter{2192}{\ensuremath{\rightarrow}}
\DeclareUnicodeCharacter{2212}{\ensuremath{-}}
\DeclareUnicodeCharacter{2265}{\ensuremath{\geq}}
\DeclareUnicodeCharacter{2264}{\ensuremath{\leq}}
\DeclareUnicodeCharacter{03B8}{\ensuremath{\theta}}
\DeclareUnicodeCharacter{0394}{\ensuremath{\Delta}}
\DeclareUnicodeCharacter{03B4}{\ensuremath{\delta}}
% Let long code identifiers break after an underscore, and loosen spacing.
\newcommand{\ttbreakunderscore}{\textunderscore\allowbreak}
\let\origtexttt\texttt
\renewcommand{\texttt}[1]{\origtexttt{\let\_\ttbreakunderscore #1}}
\sloppy
\usepackage{bookmark}
\IfFileExists{xurl.sty}{\usepackage{xurl}}{} % add URL line breaks if available
\urlstyle{same}
\hypersetup{
  colorlinks=true,
  linkcolor={Maroon},
  filecolor={Maroon},
  citecolor={Blue},
  urlcolor={Blue},
  pdfcreator={LaTeX via pandoc}}

\title{Improvement Pathways}
\author{}
\date{}

\begin{document}
\maketitle

{
\setcounter{tocdepth}{3}
\tableofcontents
}
\section{Improvement Pathways: Specification, Proofs, and Acceptance
Criteria}\label{improvement-pathways-specification-proofs-and-acceptance-criteria}

\textbf{Revision:} 2026-09-29 (research fork). Supersedes the 2026-09-11
revision and every status block added to it up to 2026-09-23. The
narrative of what happened is
\texttt{PROJECT\_HISTORY\_AND\_SYSTEM\_DESCRIPTION.md}; dated decisions
and verdicts are \texttt{docs/PREREGISTRATION.md}. "The audit" below is
the 2026-09-23 audit this revision acts on, removed from the tree on
2026-10-01; read it with
\texttt{git\ show\ cb45743:docs/AUDIT\_2026-09-23\_D12\_AND\_PATHWAY\_A.md}.

\textbf{Target venue:} not yet chosen for this fork.

\textbf{Scope.} \emph{Programme 1} (this revision) is a per-level
reinforcement-learning compaction-trigger controller for leveled
RocksDB. It minimises a priced cost of write, read and space
amplification under a chosen priority (reads, writes or space), acting
only through RocksDB\textquotesingle s own compaction scores, so that
RocksDB keeps sole authority over which files are compacted.
\emph{Programme 2} (deferred) is SLO-constrained operation: the
constrained objective, the guard (Pathway E) and phase-aware objective
switching (Pathway F).

\textbf{Status.} Specification. No arm has run under this revision.
Every result below is either proved from its stated assumptions or
explicitly labelled a conjecture, an approximation or a measured fact
with its source. Results of the 2026-09-11 programme are summarised in
history §18.1 and are not re-scored.

\textbf{Reading order.} §0 (what changed) → §1 (notation and
assumptions) → Pathway D (objective and cost model) → Pathway A
(actuation) → Pathway G (propagation across levels) → Pathway H (RL
architecture) → Pathway B (workloads) → Pathway C (comparator) →
Pathways E and F (Programme 2) → execution order → global acceptance →
references.

\textbf{Referencing convention.} Numbered results use a dot (Lemma D.7,
Theorem A.1). Sections inside a pathway use a section sign (D §3, H §5).
Acceptance criteria use a prefix per pathway: retained criteria keep
their 2026-09-11 names (B-1, C-1, C-2, C-6, E-\emph{, F-}); new criteria
use new prefixes (OBJ, ACT, PROP, ARCH, WL, CMP) so that no new
criterion reuses the name of one already scored in
\texttt{docs/PREREGISTRATION.md}.

\textbf{RocksDB version.} Upstream base
\texttt{7ea2d73655025332855864f0b8d6d2dbdad3f336} (contract v3,
\texttt{rocksdb\_base}). The contract\textquotesingle s last recorded
fork revision is \texttt{6ad9f6b79} (parent \texttt{71627a1cd}); the
root branch records \texttt{25468bbaa}, three fork commits later.
(\texttt{81d2742bf}, pinned 2026-09-12, is contract v2\textquotesingle s
revision.) Every RocksDB citation in this revision is to
\texttt{25468bbaa}, read with
\texttt{git\ -C\ lib/rocksdb\ show\ 25468bbaa:\textless{}path\textgreater{}},
since the submodule\textquotesingle s working tree can lag the recorded
commit. Programme 1 adds one option (Pathway A), so its binary differs
from every earlier one and every static measurement is regenerated on it
(Pathway C). Every cited line number and option semantic must be
re-verified against the commit that binary is built from.

\textbf{Build and I/O.} Unchanged from 2026-09-11:
\texttt{-march=znver5} through RocksDB\textquotesingle s
\texttt{PORTABLE} variable, binary SHA-256 inside
\texttt{experiment\_fingerprint}. Direct I/O is pinned \textbf{off} for
development (a 21× throughput penalty was measured) and fingerprinted as
\texttt{dio0}. Paper runs may use \texttt{dio1}; they then form their
own comparator chain and are never pooled with \texttt{dio0} runs.
Latency figures under \texttt{dio0} are warm-cache and are reported as
diagnostics only.

\begin{center}\rule{0.5\linewidth}{0.5pt}\end{center}

\subsection{0. What changed in this revision, and
why}\label{0-what-changed-in-this-revision-and-why}

\subsubsection{§0.1 The premise behind Pathway A was an artifact of the
controller}\label{01-the-premise-behind-pathway-a-was-an-artifact-of-the-controller}

The 2026-09-11 revision built Pathway A (capacity co-design) so that the
controller could defer compaction "without the depth cascade" measured
in history §10.7, where learned arms deepened the tree by one to three
levels and paid for it in reads. That depth growth came from the
controller that produced it, not from deferral as a mechanism.

\begin{itemize}
\tightlist
\item
  \textbf{The learning setup was never re-derived when the objective
  changed.} The project\textquotesingle s goal moved from latency and
  stall reduction, to a strict three-way improvement of write, read and
  space amplification, to "reduce reads with at most 2\% more writes".
  The reward, state and horizon were carried forward through both
  changes. The reward priced write amplification, probes, scan work and
  latency, and treated space as a quantity to minimise (history §6.3).
  The state had eleven stall-era pressure inputs out of 31 and no
  write-amplification input, and the discount gave a horizon of about 20
  s (history §14.20, findings 1--6).
\item
  \textbf{Two different mechanisms produced the depth growth, and
  neither was deferral as such.} On the uniform workload (2026-09-03)
  the learner deferred top-of-tree compaction while being rewarded for
  lower space (history §10.7, Finding 3). On \texttt{Assoc} (D-7) it
  compacted L1 and L2 at about 40\% of target, which the action space
  allowed (history §14.20, finding 3).
\item
  \textbf{With the setup realigned (D-9 to D-12), depth did not grow.}
  Populated depth equalled the static twin\textquotesingle s in all 18
  D-12 arms (audit §3).
\item
  \textbf{Pathway A\textquotesingle s own measurement already said depth
  was a minor write cost.} Depth accounted for at most 35\% of the
  learner\textquotesingle s extra writing (A-0, history §18.1).
\end{itemize}

\textbf{Consequence.} Pathway A\textquotesingle s motivation is
withdrawn. Its mechanism --- changing a level\textquotesingle s
effective capacity through its compaction score --- survives as one of
the controller\textquotesingle s actions (Pathway A), with its theory
corrected.

\subsubsection{§0.2 What D-12 and the audit
established}\label{02-what-d-12-and-the-audit-established}

\begin{itemize}
\tightlist
\item
  \textbf{The first run with correct instruments found no read gain at
  write parity} on \texttt{Assoc} at 10M. At T=2 all of the extra
  writing (+0.68 GB per run) fell in the first 30 s after the controller
  resumed, before the first gradient step (38--42 s), while random
  exploration deferred due levels during the post-load backlog. After 30
  s the learner wrote 0.07 GB \emph{less} than native. Part of the T=2
  deficit was session drift: the same static setting writes 2.9--4.1\%
  more in later sessions.
\item
  \textbf{The constrained objective had almost no room at its
  comparators.} At T=2 and T=10 the comparator was already the
  lowest-read static setting within +2\% write, and at T=10 the
  lowest-read of all 38 static settings (audit §2).
\item
  \textbf{Stale answers were real but not the cause.} 27--39\% of
  controller answers were rejected over a run because the tree had
  changed since the snapshot. The prior-only arm makes the same round
  trip and matched native at T=10 (audit §6).
\item
  \textbf{Several instruments carried the post-load transient:} the
  write and scan multipliers and the guard\textquotesingle s limits. The
  guard\textquotesingle s conditional override rate was 6--17 times its
  limit over the whole run and 0--1.4\% after the first 30 s.
\item
  \textbf{The Double-DQN target ignored the next state\textquotesingle s
  action mask} (\texttt{agent.py} \texttt{train\_step}); its effect was
  never measured.
\item
  \textbf{Static capacity expansion never lowered write cost} in 27 runs
  (merges out of L0 grew 34--74\%), and lowered read cost only by
  removing a level.
\end{itemize}

\subsubsection{§0.3 The objective
changes}\label{03-the-objective-changes}

The constrained objective --- minimise point-read amplification subject
to write parity at a 2\% margin and to space, latency and stall bounds
--- is replaced by a \textbf{priced cost with a priority mode} (Pathway
D). The controller trades the three amplifications at stated prices and
weights the prioritised one by a factor \(\beta^\star\). There are no
hard limits. Constraints, the guard and latency bounds move to Programme
2.

\subsubsection{§0.4 Fixes adopted from the parallel
stream}\label{04-fixes-adopted-from-the-parallel-stream}

\begin{itemize}
\tightlist
\item
  \textbf{Measure only a settled tree.} After the load, no operation is
  issued until compaction has caught up and the tree has stayed settled
  for a hold window; the measured phase and the controller both start at
  the first \texttt{mixgraph} operation (H §5). This removes the
  start-up exploration damage of D-12, and every \texttt{mixgraph}
  operation is scored.
\item
  \textbf{In-process inference.} Inference runs in C++ inside the
  RocksDB process; Python only trains, and pushes weights at a fixed
  period (H §6). No answer can be stale by more than one decision.
\end{itemize}

\subsubsection{§0.5 Corrections to the 2026-09-11
revision}\label{05-corrections-to-the-2026-09-11-revision}

{\def\LTcaptype{none} % do not increment counter
\begin{longtable}[]{@{}
  >{\raggedright\arraybackslash}p{(\linewidth - 6\tabcolsep) * \real{0.0132}}
  >{\raggedright\arraybackslash}p{(\linewidth - 6\tabcolsep) * \real{0.3946}}
  >{\raggedright\arraybackslash}p{(\linewidth - 6\tabcolsep) * \real{0.3946}}
  >{\raggedright\arraybackslash}p{(\linewidth - 6\tabcolsep) * \real{0.1776}}@{}}
\toprule\noalign{}
\begin{minipage}[b]{\linewidth}\raggedright
\#
\end{minipage} & \begin{minipage}[b]{\linewidth}\raggedright
2026-09-11 statement
\end{minipage} & \begin{minipage}[b]{\linewidth}\raggedright
Correction
\end{minipage} & \begin{minipage}[b]{\linewidth}\raggedright
Where
\end{minipage} \\
\midrule\noalign{}
\endhead
\bottomrule\noalign{}
\endlastfoot
1 & Pathway A lets the controller "defer without the depth cascade" of
§10.7 & The cascade was produced by the misaligned controller and was
absent in D-12 & §0.1 \\
2 & Theorem A.1: a released burst is \(\kappa C_i\) & The burst is
\((\kappa-1)C_i\) at \(m_i=1\); in general \((\varphi_i - m_i)^+C_i\) &
Theorem A.1 \\
3 & Theorem A.2(iii): the survival-weighted optimal profile & Needs some
multipliers below 1, which \(s \ge 1\) could not express; multipliers
below 1 are now allowed. The profile is static & Theorem A.2 \\
4 & Corollary A.3: removing a level needs \(s = T\) & Premise failed:
\(s=2\) removed a level at T=10. Superseded: depth cannot fall during a
run, so removal is static & Lemma A.6, Corollary A.7 \\
5 & Corollary A.4 and P0-7: write-accounting convention M3 & Replaced by
an exact accounting identity whose one modelled quantity, the overlap
constant, is measured & Lemmas D.7, D.8 \\
6 & Theorem B.1: elision ceiling on \(W-1\) & The proof evaluates a
constant-survival formula at a pooled floor, which the floor does not
justify. In the stage model without trivial moves the saving is limited
by the garbage available: about 26--32\% of whole-run \(W-1\) on
\texttt{Assoc} at T = 2, about half the formula\textquotesingle s 58\%.
In the measured phase the budget is weak, native already drops most of
it, and where drops happen binds & Propositions B.1′, B.1″ \\
7 & B.1 and B.2 tables, \(\kappa\) values & Measured on the uniform
workload with the defective live-data estimate (D-3); replaced by
measured values & Pathway B \\
8 & §B and §D: cumulative garbage "is large"; §C: "almost no resident
garbage" & Resolved by D-3: on \texttt{Assoc}, resident \(S\) is
1.03--1.09 and \(g_{\text{flow}} = 0.278\) & Pathway B \\
9 & Commentary on Proposition B.3: a state-dependent policy "can
strictly beat every member of \(\Theta\)" & Unproven; no learned or
hand-written policy on record lies below the static frontier &
Conjecture B.3′ \\
10 & Gate 3b runs on the uniform workload & Contradicted D-1 (every gate
on \texttt{Assoc}) & Execution order \\
11 & D-4: \(\lambda_W\) diverges where the write constraint is
infeasible & Could never apply, since native RocksDB always meets its
own bound; retired with the constrained objective & history §18.1 \\
12 & Gate costs & 5--7 times too high (audit §4); new gates are priced
by run count once run length is fixed & Execution order \\
13 & D-12 verdict: the T=2 extra writing is a learned deferral lever &
Start-up exploration during the backlog plus session drift & history
§18.1 \\
14 & Proposition F.2 bounds the anticipation term by the read-phase cost
over the first \(\max(0, \Delta_{\text{shape}} - \ell)\) & That quantity
is the loss a predictor still leaves, not the anticipation term;
restated with proof & Proposition F.2 \\
\end{longtable}
}

\subsubsection{§0.6 Dated decisions this revision requires before any
run}\label{06-dated-decisions-this-revision-requires-before-any-run}

These belong in \texttt{docs/PREREGISTRATION.md}, each dated and
committed before the run it governs. This document states the theory; it
does not substitute for the dated record.

\begin{enumerate}
\def\labelenumi{\arabic{enumi}.}
\tightlist
\item
  The objective change (§0.3), the headline priority factor
  \(\beta^\star\), the reference operation rate \(\bar q\) per workload,
  and the two money prices the others are converted with: the instance
  price per device-second (one core-second, the instance price over its
  cores; PREREGISTRATION D-15) and the storage price \(c_s > 0\) per
  byte-second (D §1).
\item
  An amendment to P0-7 (write accounting): exact identity plus measured
  overlap constant.
\item
  The actuation option and its bounds (Pathway A).
\item
  The workload suite and its fingerprint fields (Pathway B).
\item
  The static comparator class \(\Theta_s\) (Pathway C).
\item
  The settle rule and its hold window \(h_w\) (H §5).
\item
  Run length, and the admission test\textquotesingle s fixed parameters:
  reference level, statistics, margins \(\delta_{\text{adm},s}\) and
  \(\omega_{\max}\), block length \(b\), minimum turnovers per level
  \(n_{\min}\) (Gate N1), and PROP-1b\textquotesingle s one-step model
  and margin \(\delta_{\text{kern}}\).
\item
  Repeat counts per cell (C §3).
\item
  Retirement, for this fork, of P0-1 as a constraint (scans are now
  priced), P0-3, P0-6, P1-15, P1-16, P1c-22 and P1c-23; P0-4 (latency)
  moves to Programme 2.
\item
  The claim the paper will make (Global acceptance), recorded before
  Gate N5.
\item
  Whether the neighbour charge of H §3 uses the
  neighbour\textquotesingle s full value or a head fitted to its
  attributed cost alone (the echo).
\item
  The stall rule\textquotesingle s margins \(\delta_{\text{stall}}\) and
  \(\delta_{\text{thr}}\) (Global acceptance).
\end{enumerate}

\begin{center}\rule{0.5\linewidth}{0.5pt}\end{center}

\subsection{1. Notation and standing
assumptions}\label{1-notation-and-standing-assumptions}

\subsubsection{§1.1 Notation}\label{11-notation}

{\def\LTcaptype{none} % do not increment counter
\begin{longtable}[]{@{}
  >{\raggedright\arraybackslash}p{(\linewidth - 4\tabcolsep) * \real{0.3652}}
  >{\raggedright\arraybackslash}p{(\linewidth - 4\tabcolsep) * \real{0.3652}}
  >{\raggedright\arraybackslash}p{(\linewidth - 4\tabcolsep) * \real{0.2496}}@{}}
\toprule\noalign{}
\begin{minipage}[b]{\linewidth}\raggedright
Symbol
\end{minipage} & \begin{minipage}[b]{\linewidth}\raggedright
Meaning
\end{minipage} & \begin{minipage}[b]{\linewidth}\raggedright
Source
\end{minipage} \\
\midrule\noalign{}
\endhead
\bottomrule\noalign{}
\endlastfoot
\(L\) & index of the last populated level; levels are \(0..L\) &
\texttt{levelstats} \\
\(T\) & size ratio between adjacent level targets &
\texttt{max\_bytes\_for\_level\_multiplier} \\
\(C_i = C_1T^{i-1}\) & nominal target of level \(i \ge 1\); \(C_1\) =
\texttt{max\_bytes\_for\_level\_base} & options \\
\(m_i\) & target multiplier of level \(i\); \(m_0 \equiv 1\) &
\texttt{level\_target\_multipliers} (Pathway A) \\
\(K_0\) & L0 compaction trigger (file count) &
\texttt{level0\_file\_num\_compaction\_trigger} \\
\(k_0(t)\), \(F\) & L0 file count; flush file size & telemetry \\
\(F_{\text{sst}}\) & SST target file size
(\texttt{target\_file\_size\_base}, 512 KiB here) & options \\
\(B_i\) & bytes at level \(i\) not being compacted & telemetry \\
\(\varphi_i = B_i/C_i\) & fill of level \(i\) against its nominal target
& derived \\
\(s_i = \varphi_i/m_i\) & compaction score of level \(i \ge 1\) &
RocksDB \\
\(S, O, X\) & per compaction: source bytes, overlap bytes read from the
level below, output bytes & event log \\
\(\rho_i = (X-O)/S\) & pass-through of level \(i\): net bytes landing
below per source byte of a merge, summed over merges with trivial moves
excluded; 1 if nothing is dropped & event log \\
\(o_i = O/S\) & overlap ratio of level \(i\) & event log \\
\(d = S + O - X = (1-\rho_i)S\) & bytes dropped by one compaction &
event log \\
\(\lambda_i\), \(\mu_i\) & inflow to level \(i\) (net bytes added by
compactions from above) and outflow (source bytes leaving), bytes/s &
event log \\
\(\bar\lambda_i\) & mean inflow over a window & derived \\
\(\tau_i = C_i/\bar\lambda_i\) & turnover time: time for the inflow to
deliver one nominal level of bytes & derived \\
\(\theta_i = t/\tau_i\) & the level\textquotesingle s own clock, in
turnovers & derived \\
\(a_i\) & source bytes leaving level \(i\) per user byte, by merge or by
trivial move (\(a_0\): bytes leaving L0) & derived \\
\(a_F\) & flush bytes per user byte; \(a_F = a_0\) unless an intra-L0
compaction drops entries & derived \\
\(t_i\), \(\xi_i = t_i/a_i\) & bytes leaving level \(i\) by trivial move
(a file relinked into level \(i+1\) without being rewritten) per user
byte, and their share of \(a_i\) & event log \\
\(\tilde\rho_i = \xi_i + (1-\xi_i)\rho_i\) & net pass-through of level
\(i\), trivial moves included; \(= a_{i+1}/a_i\) in steady state &
derived \\
\(\pi_i = \prod_{k=1}^{i-1}\tilde\rho_k\) & share of
L1\textquotesingle s inflow that reaches level \(i\); \(= a_i/a_1\) in
steady state & derived \\
\(f_i\) & effective fanout: \(f_0 = m_1C_1/(K_0F)\);
\(f_i = Tm_{i+1}/m_i\) for \(1 \le i \le L-2\);
\(f_{L-1} = B_L/(m_{L-1}C_{L-1})\) & derived \\
\(c_i = o_i/f_i\) & overlap constant of level \(i\) & measured \\
\(W\) & write amplification: (flush + compaction bytes written) / user
bytes written & evaluator \\
\(R_f\) & filter probes per Get (the 2026-09-11 "point-read
amplification") & existing counter \\
\(R_{blk}\) & probes per Get that pass the filter and read a data block
& Pathway D \\
\(R_{sk}\) & sorted runs seeked per scan (P1c-20) & existing counter \\
\(S\) (metric) & settled SST bytes / garbage-free bytes (D-3) &
evaluator \\
\(H(t)\) & SST bytes of the current version
(\texttt{rocksdb.live-sst-files-size}): a running
compaction\textquotesingle s inputs count until it installs and its
outputs from then on; obsolete files still pinned by older versions or
awaiting deletion do not count & telemetry \\
\(u\), \(q_{pt}\), \(q_{sc}\) & user bytes written, Gets and scans, per
second & telemetry \\
\(q\), \(\bar q\) & all operations per second; a reference operation
rate, fixed in advance per workload and recorded with the prices
(OBJ-2), since it sets the weight of space against the other terms &
telemetry; preregistration \\
\(N_i = C_i\,q/\bar\lambda_i\) & operations served per turnover of level
\(i\) (\(q\) and \(\bar\lambda_i\) over the same window) & derived \\
\(c_w, c_f, c_{blk}, c_{sk}, c_s\) & prices: per byte written, per
filter probe, per block-reading probe, per run seek, per byte held per
second & measured once per machine \\
\(\beta = (\beta_W, \beta_R, \beta_S)\) & priority weights;
\(\beta^\star > 1\) is the factor on the prioritised term & run
configuration \\
\(J_\beta(\pi)\) & priced, priority-weighted cost of policy \(\pi\) over
the measured phase & Pathway D \\
\(\mathcal C_W, \mathcal C_R, \mathcal C_S\) & unweighted priced run
costs of writes, reads (three terms together) and space;
\(J_\beta = \beta_W\mathcal C_W + \beta_R\mathcal C_R + \beta_S\mathcal C_S\)
& Pathway D \\
\(\Theta_s\) & the static configuration class & Pathway C \\
\(\theta^\star_\beta\) & the static configuration minimising \(J_\beta\)
on a workload & Pathway C \\
\(S_{\text{flow}} = N_{\text{written}}/N_{\text{live}}\),
\(g_{\text{flow}} = 1 - 1/S_{\text{flow}}\) & flow space ratio and
cumulative garbage fraction & as 2026-09-11, measured denominator
(D-3) \\
\end{longtable}
}

\subsubsection{§1.2 Standing assumptions}\label{12-standing-assumptions}

\begin{itemize}
\tightlist
\item
  \textbf{A1 (static ladder).} Leveled compaction with
  \texttt{level\_compaction\_dynamic\_level\_bytes\ =\ false}; the
  nominal targets \(C_i\) are fixed for the run. The multiplier option
  is rejected otherwise (Pathway A). This deviates from the RocksDB
  default since v8.4 and is declared in the paper.
\item
  \textbf{A2 (one run per level).} Levels \(\ge 1\) each hold one sorted
  run, so one version per key per level; L0 holds \(k_0\) overlapping
  runs. Exceptions to state in the paper: intra-L0 compaction,
  subcompactions, trivial moves, SST ingestion.
\item
  \textbf{A3 (eligibility).} Level \(i \ge 1\) is eligible for
  compaction iff \(s_i = B_i/(m_iC_i) \ge 1\). RocksDB ranks eligible
  levels by score, and admission depends on free compaction slots.
  Results use eligibility only unless they say otherwise.
\item
  \textbf{A3′ (L0).} L0\textquotesingle s score is
  \(\max(k_0/K_0,\; \text{L0 bytes}/C_1)\). Multipliers do not enter it.
  The byte branch caps the effective trigger at about \(C_1/F\).
\item
  \textbf{A4 (downward movement).} Compactions move data from level
  \(i\) to \(i+1\), or within L0 or within the last level; never upward.
\item
  \textbf{A5 (fluid approximation).} Where a result treats flows as
  continuous it says so; per-compaction granularity is noted where it
  matters.
\item
  \textbf{A6 (entries).} Unless stated otherwise: fixed entry size, no
  deletes or tombstones, and compression that does not depend on how
  entries are grouped into files.
\item
  \textbf{A7 (steady state).} A window in which each
  level\textquotesingle s mean fill is constant, so that each
  level\textquotesingle s mean inflow equals its mean outflow.
\item
  \textbf{A8 (same-phase measurement).} All compared arms are scored
  over the same operations --- every \texttt{mixgraph} operation, from
  the first (\(n_w\)) to the last, then the drain --- with same-session
  twins (C §3). Each arm reaches \(n_w\) on the settled tree its own
  load left (H §5); an arm not settled by then is invalid. Costs are
  totals over these operations.
\end{itemize}

\begin{center}\rule{0.5\linewidth}{0.5pt}\end{center}

\subsection{Pathway D --- Objective: a priced cost with a priority
mode}\label{pathway-d--objective-a-priced-cost-with-a-priority-mode}

\subsubsection{§0 Purpose, and the model this pathway starts
from}\label{0-purpose-and-the-model-this-pathway-starts-from}

This pathway defines what the controller minimises, how "prioritise
reads", "prioritise writes" and "prioritise space" are expressed, and
the cost model the rest of the document uses. It replaces the
constrained objective of the 2026-09-11 revision. Of that
revision\textquotesingle s Pathway D results, the shaping result is
retained as Proposition H.4; the hinge and two-timescale results are
retired with the constraints.

\textbf{The starting model.} The draft of 2026-09-29 adapted the cost
form of Dostoevsky\textquotesingle s Fluid LSM-tree {[}Dostoevsky{]}.
With \(K\) runs per upper level, \(Z\) at the last level and \(L\)
levels:

\[\text{RA} = K(L-1) + Z,\qquad \text{WA} = \frac{T}{K}(L-1) + \frac{T}{Z},\qquad \text{SA} = \frac{ZT}{T-1},\]
\[C = w\cdot\text{WA} + r\cdot\text{RA} + s\cdot\text{SA},\]

where \(w\) is the ingest rate times the cost of writing one byte, \(r\)
the lookup rate times the cost of checking one run, and \(s\) the live
data size times the cost of holding one byte for one second. The prices
turn three quantities in different units into one currency (device time
per second, or money per second {[}Cosine{]}).

\emph{Check of the plotted example.} With \(Z = K\) and \(s = 0\),
\(C = wTL/K + rKL\), minimised at \(K^\star = \sqrt{wT/r}\). At
\(T = 16\) and \(w = r = 1\) this gives \(K^\star = 4\); with
\(L \approx 3.33\), \(\text{WA} = \text{RA}
= 13.3\) and \(\text{SA} = 4\cdot16/15 = 4.27\), matching the plot.

\textbf{Five issues, and where each is fixed.}

\begin{enumerate}
\def\labelenumi{\arabic{enumi}.}
\tightlist
\item
  \textbf{The controller\textquotesingle s knobs do not change \(K\).}
  Every level \(\ge 1\) of leveled RocksDB holds one run (A2), so
  \(K = Z = 1\) whatever the multipliers are. \(K\) exists only at L0,
  where the trigger \(K_0\) plays its role. \emph{Fixed} by writing the
  model in the actual knobs: \(K_0\), the multipliers, and depth (§3).
\item
  \textbf{RA priced every run check as a random page read.} With Bloom
  filters, most checks are an in-memory filter probe. \(K(L-1)+Z\) is
  right for scans, which seek every run, and for lookups without
  filters. \emph{Fixed} by pricing filter probes, block-reading probes
  and scan seeks separately (§1, Lemma D.10).
\item
  \textbf{The write term used one convention (T per level) and the
  frozen P0-7 used another (M3).} \emph{Fixed} by an exact accounting
  identity whose only modelled quantity, the overlap constant, is
  measured (Lemmas D.7, D.8).
\item
  \textbf{SA is a worst case.} \(ZT/(T-1)\) assumes every upper-level
  entry updates a key in the last level: 2.0 at T=2, where
  \texttt{Assoc} measures 1.07--1.09. \emph{Fixed}: the evaluation uses
  measured space and the per-level reward the shadowed-garbage estimate
  of §4 (OBJ-3); the model is used only as a bound (Lemma D.14).
\item
  \textbf{"Min-max" named a weighted sum.} \emph{Fixed}: the objective
  is a priced cost with a priority factor (§2). The true min--max form
  appears only where it belongs: the Chebyshev alternative (Remark D.6)
  and suite robustness (Definition C.5).
\end{enumerate}

\subsubsection{§1 Metrics and prices}\label{1-metrics-and-prices}

\textbf{Metrics}, all over the measured phase (A8):

{\def\LTcaptype{none} % do not increment counter
\begin{longtable}[]{@{}
  >{\raggedright\arraybackslash}p{(\linewidth - 4\tabcolsep) * \real{0.0891}}
  >{\raggedright\arraybackslash}p{(\linewidth - 4\tabcolsep) * \real{0.4455}}
  >{\raggedright\arraybackslash}p{(\linewidth - 4\tabcolsep) * \real{0.4455}}@{}}
\toprule\noalign{}
\begin{minipage}[b]{\linewidth}\raggedright
Metric
\end{minipage} & \begin{minipage}[b]{\linewidth}\raggedright
Definition
\end{minipage} & \begin{minipage}[b]{\linewidth}\raggedright
Instrument
\end{minipage} \\
\midrule\noalign{}
\endhead
\bottomrule\noalign{}
\endlastfoot
\(W\) & (flush + compaction bytes written) / user bytes written. SST
output only (flush \texttt{table\_file\_creation} sizes, compaction
\texttt{total\_output\_size}): the OPTIONS file every
\texttt{SetOptions} call writes (A-Impl-5), and MANIFEST, WAL and
info-log writes, are not counted & event log, as D-11 \\
\(R_f\) & filter probes per Get: tables whose key range covers the key
and whose filter is consulted, including filter-rejected ones & existing
counter \\
\(R_{blk}\) & probes per Get that pass the filter and read a data block
(true plus false positives) &
\texttt{rocksdb.bloom.filter.full.positive}; per level through counters
keyed by the version\textquotesingle s level (§4, Gate N0) \\
\(R_{sk}\) & sorted runs seeked per scan (P1c-20) & existing counter \\
\(S\) & settled SST bytes / garbage-free bytes (D-3), at run end &
evaluator \\
\(\mathcal C_S\) & \((c_s/\bar q)\sum_n H(n)\) over the measured phase,
computed exactly as the sum, over the intervals between version
installs, of \(H\) times the operations served in the interval; \(H\) is
sampled with the operation count at every install, on every arm, native
included & telemetry, evaluator \\
\end{longtable}
}

Block-reading probes are counted logically, whether or not the block was
cached, so the metric does not depend on the page cache.

\textbf{Prices.} One currency throughout: money, as in Cosine
{[}Cosine{]}. \(c_w\) per byte written; \(c_f\) per filter probe, a CPU
cost that is not negligible on fast storage {[}Zhu et al.{]};
\(c_{blk}\) per block-reading probe; \(c_{sk}\) per run seek: each is
the device time the operation takes, measured once per machine,
converted at a fixed instance price per device-second, a device being
one core (D-15: every priced operation runs on one thread). \(c_s\) per
byte held per second is a storage price, and \(c_s > 0\) always. At
\(c_s = 0\) the space term vanishes: space priority becomes balanced
mode, and in every mode expanding levels and holding garbage cost
nothing. The two money prices are fixed in advance (§0.6) and every
price is recorded in the fingerprint (OBJ-2). Since the weight of space
against the other terms is set by a price choice, results are also
reported at \(c_s/2\) and \(2c_s\). The operation rates \(u\), \(q\),
\(q_{pt}\) and \(q_{sc}\) are measured online; the reference rate
\(\bar q\) is fixed in advance per workload (§0.6).

\textbf{Cost rate.}
\[c(t) = c_w\,\dot w(t) + q_{pt}(t)\big(c_f R_f + c_{blk}R_{blk}\big)(t) + q_{sc}(t)\,c_{sk}R_{sk}(t) + c_s H(t)\,\frac{q(t)}{\bar q},\]
where \(\dot w\) is bytes written per second. Relation to the draft:
\(w\cdot
\text{WA} = c_w u W\); \(r\cdot\text{RA}\) becomes the three read terms;
\(s\cdot\text{SA} = c_s N_{\text{live}}S = c_s H\) at the reference
rate.

\textbf{Space is charged per operation served, not per second.} Every
arm serves the same operation sequence (one client thread, a fixed seed,
\texttt{mixgraph} stopped by operation count) and is scored over the
same operations (A8), so the write and read terms integrate to counts
that do not depend on how long the run takes. A space term of \(c_sH\)
per second would not: a run that takes longer holds the same live bytes
for longer and pays for it. With the factor \(q/\bar q\),
\(\int c_sH\,(q/\bar q)\,dt = (c_s/\bar q)\sum_n H(n)\), where \(H(n)\)
is \(H\) when operation \(n\) is served. \(H\) changes only when a
version is installed, so the sum is computed exactly as the sum, over
the intervals between installs, of \(H\) times the operations served in
the interval. Space is priced by what is held while the store serves
each operation (Lemma D.15). \(c_s\) keeps its meaning --- the price of
holding one byte for one second --- when the store serves \(\bar q\)
operations per second.

\(J_\beta\) therefore contains no time. Stalls, slowdowns, foreground
throughput and the controller\textquotesingle s own overhead are not
priced in Programme 1; they are reported per arm as paired diagnostics
(OBJ-6), and pricing them belongs to Programme 2. Because they are
unpriced, a policy could lower \(J_\beta\) by deferring work until
writes stall; the stall rule (Global acceptance) bars any claim that
does. The drain serves no operations, so it carries no space cost; its
bytes written are counted. A controller arm drains under its fallback
settings (\(m \equiv 1\), \(K_0\) as configured; A-Impl-8), the
configuration it started from, and a static arm under its own
configuration, so every arm drains to the targets it would hold without
a controller and no deferred work leaves the run uncharged. Draining a
static profile at \(m \equiv 1\) instead would charge it a compaction
its configuration never performs. The rule errs against the controller:
one that ends with anchors above 1 pays for the drain, while a static
profile holding the same multipliers does not.

\textbf{Lemma D.1 (the flow form is exact).} Let the reward of an
interval \([t, t+\Delta t)\) be \(-\int_t^{t+\Delta t} c\). For any
partition of the measured phase into intervals, the rewards sum to minus
the run\textquotesingle s priced cost:
\(c_w\cdot(\text{bytes written}) + \sum(\text{read price} \times \text{count}) +
(c_s/\bar q)\int H\,q\,dt\).

\emph{Proof.} Counts and integrals are additive over disjoint intervals.
\(\blacksquare\)

\emph{Consequence.} No run-to-date ratio enters the reward, so a
start-up transient cannot inflate a multiplier the way it inflated
\(\lambda_W\) and \(\lambda_{\text{scan}}\) in D-12. The transient still
enters the cost itself, which is why the measured phase starts only
after the tree settles (H §5).

\emph{Scope.} Lemma D.1 is about the global cost. The agents of Pathway
H are trained on a different quantity: each level\textquotesingle s
attributed, priority-weighted cost divided by \(c_wC_i\), plus
counterfactual neighbour charges (H §3). Summed over levels, those
rewards do not give minus the priced cost, for four reasons. Each level
is divided by a different constant. Space is charged per level as
shadowed garbage \(g_i\) (§4), an estimate that leaves out live bytes,
not as held bytes \(H\). The block read of the table where a Get finds
its key goes to a shared bucket no agent is rewarded on (§4). And a
neighbour charge \(X_{i+1}\) is a signed estimate of how level
\(i\)\textquotesingle s action, against holding, changes level
\(i+1\)\textquotesingle s future cost; level \(i+1\) pays the realised
change again in its own attributed cost, so the sum carries each
action\textquotesingle s effect on its neighbours twice --- once
realised, once estimated --- above the cost when actions burden
neighbours and below it when they relieve them. The identity that can be
checked on a run is Proposition D.16\textquotesingle s: the attributed
write and read costs, summed over levels and over every interval of the
measured phase, plus the shared hit-read bucket, before normalisation
and without neighbour charges, equal the global write and read costs
(OBJ-1, OBJ-4). The agents\textquotesingle{} discounted per-level
returns are a training surrogate: a joint policy that every
agent\textquotesingle s return rates optimal need not minimise
\(J_\beta\) (H §3), and every run is judged on \(J_\beta\).

\subsubsection{§2 Priority modes}\label{2-priority-modes}

\textbf{Definition (priced cost with priority).} For weights
\(\beta = (\beta_W,
\beta_R, \beta_S) > 0\):
\[C_\beta(t) = \beta_W\,c_w\dot w + \beta_R\big[q_{pt}(c_fR_f + c_{blk}R_{blk}) + q_{sc}c_{sk}R_{sk}\big] + \beta_S\,c_sH\,q/\bar q,\]
and \(J_\beta(\pi) = \mathbb{E}\int C_\beta\,dt\) over the measured
phase, the expectation over seeds. Write \(\mathcal C_W(\pi)\),
\(\mathcal C_R(\pi)\), \(\mathcal C_S(\pi)\) for the three unweighted
priced run costs (write, the three read terms together, space), so
\(J_\beta = \beta_W \mathcal C_W + \beta_R \mathcal C_R + \beta_S \mathcal C_S\).

\textbf{Modes.} \(\beta^\star > 1\) is the priority factor.

{\def\LTcaptype{none} % do not increment counter
\begin{longtable}[]{@{}
  >{\raggedright\arraybackslash}p{(\linewidth - 2\tabcolsep) * \real{0.4900}}
  >{\raggedright\arraybackslash}p{(\linewidth - 2\tabcolsep) * \real{0.4900}}@{}}
\toprule\noalign{}
\begin{minipage}[b]{\linewidth}\raggedright
Mode
\end{minipage} & \begin{minipage}[b]{\linewidth}\raggedright
\((\beta_W, \beta_R, \beta_S)\)
\end{minipage} \\
\midrule\noalign{}
\endhead
\bottomrule\noalign{}
\endlastfoot
balanced (pure priced cost) & \((1, 1, 1)\) \\
read priority & \((1, \beta^\star, 1)\) \\
write priority & \((\beta^\star, 1, 1)\) \\
space priority & \((1, 1, \beta^\star)\) \\
\end{longtable}
}

\textbf{Proposition D.2 (gain--loss form).} Fix any reference policy
\(\pi_0\), for example native RocksDB with the same options. For a mode
with prioritised term \(P \in \{W, R, S\}\), define
\[G_P(\pi) = \beta^\star\big[\mathcal C_P(\pi_0) - \mathcal C_P(\pi)\big] - \sum_{X \ne P}\big[\mathcal C_X(\pi) - \mathcal C_X(\pi_0)\big].\]
Then \(J_\beta(\pi) = J_\beta(\pi_0) - G_P(\pi)\). Minimising
\(J_\beta\) is therefore the same as maximising "\(\beta^\star\) times
the decrease of the prioritised cost, minus the increases of the other
two costs", and the reference cancels out of the choice.

\emph{Proof.} Expand
\(J_\beta(\pi_0) - J_\beta(\pi) = \beta^\star(\mathcal C_P(\pi_0) -
\mathcal C_P(\pi)) + \sum_{X\ne P}(\mathcal C_X(\pi_0) - \mathcal C_X(\pi))\).
\(\blacksquare\)

\emph{The three modes written out.} "Increase" is signed: if a
non-prioritised cost falls, that counts in the policy\textquotesingle s
favour at its price.

\begin{itemize}
\tightlist
\item
  \textbf{Read priority:} maximise \(\beta^\star\cdot\)(read-cost
  decrease) \(-\) (write-cost increase) \(-\) (space-cost increase).
\item
  \textbf{Write priority:} maximise \(\beta^\star\cdot\)(write-cost
  decrease) \(-\) (read-cost increase) \(-\) (space-cost increase).
\item
  \textbf{Space priority:} maximise \(\beta^\star\cdot\)(space-cost
  decrease) \(-\) (read-cost increase) \(-\) (write-cost increase).
\end{itemize}

No reference policy is needed online: by Lemma D.1 and Proposition D.2,
a global per-interval reward \(-\int C_\beta\) would implement all three
modes. The agents of Pathway H receive per-level rewards built from it
instead (H §3): attributed write and read costs, a garbage charge for
space, and neighbour charges. \(\beta\) enters every one, so the mode
reaches every agent, but the per-level rewards implement it only as far
as they approximate the difference reward (H §3, "The approximation").
Runs are scored on \(J_\beta\).

\textbf{Proposition D.3 (exchange rate).} Suppose that near an optimum
the achievable operating points form a smooth curve along which only the
prioritised cost \(\mathcal C_P\) and one other cost \(\mathcal C_X\)
change. At an interior minimiser of \(J_\beta\),
\(-d\mathcal C_P/d\mathcal C_X = 1/\beta^\star\). In metric units,
\(-dP/dX = \text{price}_X / (\beta^\star\,\text{price}_P)\).

\emph{Proof.} First-order condition along the curve:
\(\beta^\star\,d\mathcal C_P + d\mathcal C_X = 0\). \(\blacksquare\)

In plain terms, the controller accepts one extra unit of another cost
only if it buys at least \(1/\beta^\star\) units of reduction in the
prioritised cost. \(\beta^\star\) says how much more the prioritised
metric is worth than its price.

\textbf{Proposition D.4 (a large enough factor gives strict priority on
a finite class).} Let \(A\) be a finite set of operating points, such as
the static class. There is a finite \(\bar\beta\) such that for every
\(\beta^\star > \bar\beta\), every minimiser of \(J_\beta\) over \(A\)
minimises \(\mathcal C_P\) over \(A\) and, among those points, minimises
\(\sum_{X\ne P}\mathcal C_X\).

\emph{Proof.} If every point has the same \(\mathcal C_P\) the claim is
immediate. Otherwise let \(P_{\min} = \min_A \mathcal C_P\),
\(\Delta = \min\{\mathcal C_P(a) - P_{\min} :
\mathcal C_P(a) > P_{\min}\} > 0\) (finite set), and
\(M = \max_A\Sigma - \min_A\Sigma\) with
\(\Sigma = \sum_{X\ne P}\mathcal C_X\). Take \(\bar\beta = M/\Delta\).
If \(\mathcal C_P(a) >
P_{\min}\) then
\(J_\beta(a) \ge \beta^\star(P_{\min} + \Delta) + \min_A\Sigma >
\beta^\star P_{\min} + \max_A\Sigma \ge J_\beta(a')\) for every \(a'\)
with \(\mathcal C_P(a') = P_{\min}\). Among such \(a'\), \(J_\beta\)
differs only through \(\Sigma\). \(\blacksquare\)

\emph{Use.} A large \(\beta^\star\) reads "prioritise reads first, then
keep writes and space as low as possible"; a moderate one reads "trade
at this rate". Report \(\bar\beta\) per workload so readers know which
regime a result is in. The headline \(\beta^\star\) is fixed in advance
(§0.6), with \(\beta^\star \in \{2, 5,
10\}\) reported.

\textbf{Proposition D.5 (weighted costs select only supported points).}
A minimiser of \(J_\beta\) over a set \(A\) lies on the lower boundary
of the convex hull of \(A\)\textquotesingle s priced-cost vectors. A
Pareto-optimal point of \(A\) strictly above that boundary minimises
\(J_\beta\) for no \(\beta > 0\).

\emph{Proof.} \(J_\beta\) is linear in the priced-cost vector, and the
minimum of a linear function over \(\operatorname{conv}(A)\) is attained
at points of \(A\). If \(x\) minimises it for some \(\beta > 0\), the
hyperplane \(\{y : \langle\beta, y\rangle
= \langle\beta, x\rangle\}\) supports \(\operatorname{conv}(A)\) at
\(x\). A point strictly above the lower boundary has, for every
\(\beta\), a point of the hull with a smaller value. \(\blacksquare\)
{[}Das \& Dennis; Miettinen{]}

\emph{Consequences.} The comparator for every mode is a vertex of one
hull (Proposition C.4). A trade that lies in a dent of the frontier is
reachable by no choice of weights.

\textbf{Remark D.6 (the actual min--max form).} Minimising the worst
weighted gap to an ideal point,
\(\max_X \beta_X(\mathcal C_X - \mathcal C_X^{\text{ideal}})\) (the
weighted Chebyshev form), reaches every Pareto-optimal point for
suitable weights, including points in dents {[}Miettinen{]}. It needs
the ideal point --- the best achievable value of each cost separately
--- which a cold-started controller does not have. Programme 1 therefore
uses the weighted form. The Chebyshev form may be used offline to select
among measured configurations.

\subsubsection{§3 Cost model in the actual
knobs}\label{3-cost-model-in-the-actual-knobs}

Levels \(0..L\); knobs \(K_0\), \(m_1..m_L\), and depth. The model
explains and bounds; the attributed costs in the reward are always
measured (only the neighbour charge of H §3 uses a one-step model
prediction and learned values).

\textbf{Lemma D.7 (write accounting identity).} In steady state (A7),
with \(a_F\) the flush bytes per user byte, \(a_0\) the bytes leaving L0
per user byte (equal to \(a_F\) unless an intra-L0 compaction drops
entries), \(t_i\) the bytes leaving level \(i\) by trivial move per user
byte, \(\rho_i\) and \(o_i\) measured over merges only, and
\(a_{i+1} = (a_i - t_i)\rho_i + t_i\),
\[W = a_F + \sum_{i=0}^{L-1} (a_i - t_i)(\rho_i + o_i),\] plus, listed
separately, any intra-L0 or last-level self-compaction bytes.

\emph{Proof.} A merge sourced at level \(i\) writes
\(X = (X - O) + O = \rho_iS + O
= S(\rho_i + o_i)\) bytes. A trivial move relinks its file into level
\(i+1\) and writes nothing. In steady state the bytes leaving level
\(i\) per user byte equal those entering it, \(a_i\): \(a_i - t_i\) by
merge and \(t_i\) by trivial move. The net bytes landing in level
\(i+1\) per user byte are \((a_i - t_i)\rho_i + t_i =
a_{i+1}\). Summing flush bytes and the bytes written by every
level\textquotesingle s merges, per user byte, gives the identity.
\(\blacksquare\)

The identity is exact given measured \(\rho_i\), \(o_i\) and \(t_i\); no
convention enters. The expression for \(W\) needs no steady state: with
\(a_i - t_i\) read as the merged source bytes measured in the window, it
holds over any window, the measured phase included; only the recursion
for \(a_{i+1}\) needs A7. A RocksDB job is either a trivial move or a
merge, never part of each (\texttt{Compaction::IsTrivialMove} accepts or
rejects the whole input set), so \(t_i\) is well defined per job.
Counting trivially moved bytes as merged would overstate \(W\) by
\(\sum_i t_i(\rho_i + o_i)\). How large \(t_i\) is must be measured
(Pathway B §2 item 4), not assumed. Trivial moves happen where the level
below has gaps in its key coverage: while the tree grows, in the holes a
compaction leaves behind, and while a backlog drains, which native
RocksDB does with trivial moves (audit 2026-09-23, §1). On thirty
2026-09 uniform runs of the hand-written prior
(\texttt{unconstrained\_prior\_only}, 10M, T = 2/6/10), whole run with
the load, 42--71\% of the bytes leaving levels \(\ge 1\) moved
trivially, and charging them as merges would have overstated bytes
written by 26--70\%. Their share in the measured phase on \texttt{Assoc}
has not been measured.

\textbf{Lemma D.8 (overlap with uniformly spread keys).} If keys are
spread uniformly within each level and a compaction\textquotesingle s
source range covers a fraction \(x\) of level \(i\)\textquotesingle s
key space, then \(o_i = B_{i+1}/B_i\) at that moment. With both levels
at their effective targets, \(o_i = f_i\). For L0 → L1 with L0 files
spanning the key range, \(o_0 = B_1/(k_0F) = m_1C_1/(K_0F) = f_0\) at
the trigger.

\emph{Proof.} The source holds \(xB_i\) bytes and overlaps \(xB_{i+1}\)
bytes below. \(\blacksquare\)

\emph{The overlap constant.} Define \(c_i = o_i/f_i\), measured per
level. RocksDB\textquotesingle s \texttt{kMinOverlappingRatio} picks
files with below-average overlap {[}Sarkar et al.{]}, so \(c_i \le 1\)
is expected, and skewed keys lower it further. The conventions of the
2026-09-11 revision (M1: \(c f\) with \(c = 1\); M2: \(c = \tfrac12\);
M3: \(c(f+1)\) with \(c = \tfrac12\) {[}How to Grow{]}) differ in how
they count the source rewrite and the overlap; the identity counts both
exactly and leaves \(c_i\) to measurement. This resolves issue 3 and
amends P0-7 (§0.6).

\textbf{Corollary D.9 (garbage-free steady state).} With \(\rho_i = 1\),
\(a_F = a_0 = 1\), no trivial moves (\(t_i = 0\)), levels at their
effective targets and a common overlap constant \(c\),
\[W = 1 + L + c\sum_{i=0}^{L-1} f_i,\qquad \prod_{i=0}^{L-1} f_i = \frac{B_L}{K_0F},\]
and the product does not depend on \(m_1, \dots, m_{L-1}\).

\emph{Proof.} Lemma D.7 with \(a_i = 1\) and Lemma D.8. The product
telescopes:
\(\frac{m_1C_1}{K_0F}\cdot\prod_{i=1}^{L-2}\frac{m_{i+1}C_{i+1}}{m_iC_i}\cdot
\frac{B_L}{m_{L-1}C_{L-1}} = \frac{B_L}{K_0F}\). \(\blacksquare\)

\textbf{Lemma D.10 (read terms).} Runs are searched newest to oldest (L0
files, then L1 to L), after the memtables. A Get answered from the
memtables probes no table, so \(R_f = R_{blk} = 0\). Otherwise consider
a Get whose newest version is in the \(n\)-th table run whose key range
covers the key (a hit), or that finds no version in any table (a miss),
and let \(\varepsilon\) be the probability that a covering run which
does not hold the key passes its filter. Then:

\begin{itemize}
\tightlist
\item
  \(R_f = n\) for a hit, and the number of covering runs for a miss;
\item
  \(\mathbb{E}[R_{blk}] = \mathbf 1_{\text{hit}} + \varepsilon\,(R_f - \mathbf 1_{\text{hit}})\);
\item
  for a scan, \$R\_\{sk\} = k\_0 + \$ (number of levels \(1..L\) with a
  file at or after the seek key); a \texttt{DBIter} reseek, after more
  than \texttt{max\_sequential\_skip\_in\_iterations} hidden versions of
  one key, counts them again, which is rare here and checked with
  \texttt{rocksdb.number.reseeks.iteration}.
\end{itemize}

On \texttt{Assoc} there are no misses: the load writes every key below
\texttt{num}, and \texttt{mixgraph} draws only such keys. With a
per-level allocation such as Monkey {[}Monkey{]}, \(\varepsilon\)
becomes \(\varepsilon_j\) inside the sum.

\emph{Proof.} Filters have no false negatives, so every covering run
before the hit is probed and does not hold the key; each of those reads
a block with probability \(\varepsilon\), by linearity of expectation
(no independence between runs is needed), and the hit run reads one. A
scan seeks each L0 file and each level that has a file at or after the
seek key, once per level however many of its files the scan then enters
(P1c-20: at the recorded RocksDB commit the table iterator ticks only on
a keyed seek, and a level iterator moving to its next file does not tick
again). \(\blacksquare\)

\textbf{Proposition D.11 (static optimum of the L0 trigger).} Assume no
garbage dropped at L0 (\(\rho_0 = 1\)), uniform-key overlap at L0 (Lemma
D.8; each L0 file then overlaps L1, so none moves trivially and
\(t_0 = 0\)), a time-average L0 file count \(\bar k_0 = K_0/2 +
\delta_0\) with \(\delta_0\) not depending on \(K_0\) (it accounts for
flushes that arrive during an L0 compaction, and while the single
compaction slot runs another level\textquotesingle s job; independence
from \(K_0\) is an assumption, checked on the static sweep of Gate N2),
Gets and scans that probe every L0 run, and fixed prices and rates. The
part of \(C_\beta\) that depends on \(K_0\) is
\[\frac{\beta_W c_w u\,m_1C_1}{K_0F} + \beta_R\frac{K_0}{2}\big[q_{pt}(c_f + \varepsilon c_{blk}) + q_{sc}c_{sk}\big],\]
minimised over real \(K_0 > 0\) at
\[K_0^\star = \sqrt{\frac{2\,\beta_W c_w u\,(m_1C_1/F)}{\beta_R\big[q_{pt}(c_f + \varepsilon c_{blk}) + q_{sc}c_{sk}\big]}}.\]
Admissible triggers are the integers in \([2, K_{\text{cap}}]\) with
\(K_{\text{cap}}
= \min(\lfloor C_1/F\rfloor, K_{\text{slow}} - 1)\) (A3′), and the best
admissible trigger is one of the two integers next to \(K_0^\star\),
clamped to that range.

\emph{Proof.} By Lemma D.7 only the overlap term
\(a_0o_0 = m_1C_1/(K_0F)\) of the L0-sourced writes depends on \(K_0\);
by Lemma D.10 each L0 run adds one filter probe and \(\varepsilon\)
block reads per Get and one seek per scan. So the cost is
\(g(K) = A/K + BK\) with \(A, B > 0\), strictly convex on \(K > 0\) with
minimiser \(\sqrt{A/B}\). Strict convexity makes the best integer one of
the two neighbours of the real minimiser. \(\blacksquare\)

This has the same shape as the draft\textquotesingle s
\(K^\star = \sqrt{wT/r}\), with the L1-to-flush fanout \(m_1C_1/F\) in
place of \(T\). \(K_0^\star \propto \sqrt{\beta_W/\beta_R}\): read
priority lowers the trigger and write priority raises it.

\textbf{Corollary D.12 (a provable adaptivity gap for the L0 trigger).}
Suppose the workload alternates between phases whose rates
\((u, q_{pt}, q_{sc})\) give different best admissible triggers under
Proposition D.11. Then any fixed trigger costs strictly more than the
trigger switched per phase, not counting the cost of each switch.

\emph{Proof.} Each phase\textquotesingle s cost is strictly convex in
\(K_0\) with its own minimiser. A fixed \(K_0\) is at most one
phase\textquotesingle s minimiser, so it is strictly worse in at least
one phase and no better in any. \(\blacksquare\)

This predicts, in advance, a positive phase-adaptivity gap
\(\mathcal G\) (Pathway B) for one knob. A real controller also pays the
switching transient: the L0 contents at the switch.

\textbf{Proposition D.13 (interior multipliers and depth in steady
state).} Under Corollary D.9\textquotesingle s assumptions, with \(L\)
and \(K_0\) fixed:

\begin{itemize}
\tightlist
\item
  (i) \(W\) is minimised over the interior multipliers exactly when all
  fanouts are equal, \(f_i = (B_L/(K_0F))^{1/L}\). That profile is
  static.
\item
  (ii) (This part does not use Corollary D.9\textquotesingle s
  assumptions.) Up to the L0 run count, \(R_{sk}\) does not depend on
  the interior multipliers unless a level empties, and \(R_f\) and
  \(R_{blk}\) depend on them only through where Gets find their key.
  Raising \(m_i\) delays every version\textquotesingle s descent below
  level \(i\), so a Get whose key\textquotesingle s newest version would
  already have crossed a boundary at or below level \(i\) finds it one
  level higher, saving one filter probe (and \(\varepsilon\) block
  reads) per level it no longer passes. Under uniform access the moved
  share is small for the upper levels, which hold little of the data;
  for the levels just above the last it is the same shift that, taken
  far enough, empties the last level (audit §3). Under skew it can be
  material at the upper levels too, but only for Gets whose key is
  rewritten during the run. In db\_bench\textquotesingle s
  \texttt{mixgraph} the key and the operation type come from one random
  draw, so on \texttt{Assoc} about 44\% of Gets read keys that no Put
  writes during the measured phase, and on the uniform control no Get
  ever does. The moved share is measured from per-level hit shares with
  Gate N0\textquotesingle s counters (§4), not assumed. The interior
  multipliers also reach the L0 run count through \(\delta_0\)
  (Proposition D.11): \(m_1\) directly, since it sets the size, and so
  the duration, of every L0→L1 merge, and every \(m_i\) through the
  single compaction slot, whose busy time depends on each
  level\textquotesingle s merge sizes. Flushes that arrive meanwhile
  raise the mean L0 run count, adding probes to every Get that reaches
  the tables and seeks to every scan; merges out of L0 grew 34--74\% at
  \(s = 1.5\)--2 (audit §3). For a fixed profile all of this is static.
\item
  (iii) The space bound of Lemma D.14 increases in every \(m_i\).
\item
  (iv) Treat \(L\) as continuous with equal fanouts
  \(f = (B_L/(K_0F))^{1/L}\). The \(C_\beta\)-optimal fanout solves
  \[c\,f(\ln f - 1) = 1 + b/A,\] where \(A = \beta_Wc_wu\) and
  \(b = \beta_R[q_{pt}(c_f + \varepsilon c_{blk})h +
  q_{sc}c_{sk}]\) is the read cost rate one more level adds, with \(h\)
  the share of Gets that probe it. Read priority raises \(f^\star\)
  (fewer levels). At \(b = 0\) and \(c = 1\), \(f^\star \approx 3.59\).
\end{itemize}

\emph{Proof.} (i) Minimise \(\sum f_i\) with the product fixed; by the
AM--GM inequality equality of the terms is necessary and sufficient.
(ii) Lemma D.10: the read counts depend on the number of runs and the
hit position. Interior multipliers change the number of runs at levels
\(\ge 1\) only if a level empties, the hit position only by delaying
descent, and \(k_0\) only through \(\delta_0\), directly through \(m_1\)
and through the shared compaction slot. (iii) Lemma D.14. (iv) The
\(L\)-dependent cost is \(A(1 + L + cLf) + bL\). Since
\(\ln f = \ln(B_L/(K_0F))/L\), \(\frac{d}{dL}(Lf) = f(1 - \ln f)\), so
the derivative is \(A(1 + cf(1 - \ln f)) + b\); setting it to zero gives
the equation. At \(c = 1\), \(b = 0\): \(3.59\,(\ln 3.59 - 1) = 1.00\).
\(\blacksquare\)

\emph{Consequences.} Everything in (i)--(iv) is a static choice: an
interior profile, a depth fixed by base size at load time, and \(K_0\)
at fixed prices. All of it belongs to the static comparator (Corollary
C.3). Depth cannot be reduced during a run (Lemma A.6), so (iv) is a
design and comparator result only. The 2026-09-11 Corollary A.4
(\(f^\star \approx 3.59\) under M3) is the special case \(c = 1\),
\(b = 0\); the optimum moves with the measured \(c\) and with the
priority.

\textbf{Lemma D.14 (space bound).} Under A2 and A6,
\[S \le 1 + \sum_{i<L}\frac{B_i}{B_L} \le 1 + \frac{k_0F}{B_L} + \sum_{i=1}^{L-1}\frac{m_iC_i}{B_L},\]
the second inequality holding outside release transients, when
\(B_i \le m_iC_i\).

\emph{Proof.} Level \(L\) holds one version per key (A2), so the number
of distinct keys is at least the number of keys at level \(L\), and with
fixed entry size the live bytes are at least \(B_L\). Total bytes are
\(\sum_i B_i\), so \(S \le \sum_iB_i/B_L\). \(\blacksquare\)

For a full geometric ladder at \(m \equiv 1\) the bound is about
\(1 + 1/(T-1)\), which is the draft\textquotesingle s \(ZT/(T-1)\) at
\(Z = 1\). It is a worst case (every upper entry updates a last-level
key): about 2.0 at T=2 against a measured 1.07--1.09 on \texttt{Assoc}.
It is used only to bound the space cost of expansion,
\(\partial(\text{bound})/\partial m_i = C_i/B_L\).

\textbf{Lemma D.15 (live bytes do not depend on the policy).} Index the
measured phase by operation number \(n\) (A8). Suppose the arms serve
the same operation sequence, and that the bytes an entry occupies in a
table do not depend on which entries share its block or file
(A6\textquotesingle s last clause; with compression off it fails only
through per-file metadata, index and filter blocks and restart points, a
small fraction of a percent --- the garbage-free size of 108
\texttt{Assoc} runs spanned 0.006\%, D-3). Then the live bytes after
operation \(n\) are the same under every compaction policy, and with
space charged per operation served (§1) two policies\textquotesingle{}
space costs differ only through the garbage they hold:
\((c_s/\bar q)\sum_n\big(H_\pi(n) - H_{\pi_0}(n)\big) = (c_s/\bar q)\sum_n\big(G_\pi(n) - G_{\pi_0}(n)\big)\),
up to the live bytes of a memtable whose flush completes after a
different operation under different policies: at most one write buffer
at any operation (2 MiB here, with two memtables), and on average far
less, against about 3 GiB held. Entry sizes may vary; with deletes the
lemma holds with tombstones counted as garbage.

\emph{Proof.} Compaction never changes the result of a read, so the
newest version of every key after operation \(n\), and therefore the
live set, is fixed by the operations. The current
version\textquotesingle s table files (\(H\)) hold every live entry not
in the memtables, plus garbage; which entries the memtables hold after
operation \(n\) is fixed by the write sequence except for flushes still
in progress. \(\blacksquare\)

\emph{Why per operation and not per second.} Under a per-second charge
the lemma fails: two policies spend different times on the same
operations, so the difference gains
\(c_s\sum_n L(n)(\Delta t^\pi_n - \Delta t^{\pi_0}_n) \approx
c_sN_{\text{live}}(T_\pi - T_{\pi_0})\), with \(L(n)\) the live bytes
and \(\Delta t_n\) the time spent at operation \(n\). On record, runtime
moves more than space. The four base-16 MiB static triggers on
\texttt{Assoc} held within 0.6--0.8\% of one another\textquotesingle s
space (D-3), while the runtime winner led by 4.2\% at T=2 and 2.6\% at
T=10 (D-5). In the 2026-08-24 uniform matrix, learned and prior arms ran
3.6--6.3\% longer than native (history §10.6; its space figures used the
live-data estimate D-3 retired, so only the runtimes carry over). Per
second, the runtime term would be 3--7 times the space difference. A
space-priority controller would chase speed.

\emph{Consequence.} Space needs no online estimate of live data --- the
estimate D-3 retired. The evaluator prices held bytes \(H\) directly;
the per-level rewards need only the policy-dependent part, garbage, of
which §4 attributes the adjacent-level part as shadowed garbage \(g_i\)
(an estimate that undercounts; OBJ-3 checks how it tracks measured
garbage).

\subsubsection{§4 Per-level attribution}\label{4-per-level-attribution}

The agents of Pathway H each need their own share of the cost:

\begin{itemize}
\tightlist
\item
  \textbf{Writes} are charged to the source level of the compaction that
  wrote them (the job\textquotesingle s start level, not
  RocksDB\textquotesingle s per-level compaction statistics, which are
  keyed by output level); flush writes to L0. This is exact.
\item
  \textbf{Reads:} each filter probe, false-positive block read and run
  seek is charged to the level where it happens. The block read of the
  table where a Get finds its key is charged to no level. It goes to a
  shared \emph{hit-read bucket} that no agent is rewarded on. Every Get
  that reaches the tables makes exactly one such read, wherever its key
  is found (Lemma D.10), and the share of Gets answered from the
  memtables depends on the policy only through Lemma
  D.15\textquotesingle s memtable caveat. Charging it to the level that
  holds the key would charge that level for a read that merely moved up
  from a deeper level, and teach it to avoid holding recently written
  keys. This is exact, given per-level counters that take the level from
  the version being read --- \texttt{FilePicker}\textquotesingle s hit
  level in \texttt{Version::Get}, and for seeks the level at which the
  version builds the iterator (0 for each L0 file, the level
  iterator\textquotesingle s level otherwise) --- that separate the
  hit\textquotesingle s block read from false-positive ones, and that
  aggregate across threads (Gate N0). RocksDB\textquotesingle s
  per-level perf context does not qualify as it stands: it is
  thread-local, needs \texttt{perf\_level} at least
  \texttt{kEnableCount}, has no seek counter, and keys its filter
  counters by the level at which a table reader was first opened, which
  is stale after a trivial move.
\item
  \textbf{Slot blocking.} There is one compaction slot (G.4). While L0
  is due (score at least 1) but cannot compact because a job sourced at
  level \(i \ge 1\) holds the slot, the share \((k_0 - K_0)^+/k_0\) of
  L0\textquotesingle s filter probes, false-positive block reads and
  seeks --- the files beyond its trigger --- is charged to level \(i\)
  instead of to L0, for every operation served in that span. The rule
  moves cost between levels and leaves the total unchanged. It is an
  attribution rule, not an exact counterfactual: had the slot been free,
  flushes arriving during L0\textquotesingle s own merge would still
  have added files. It is the channel through which an interior level
  changes read cost during a run (Proposition D.13(ii)\textquotesingle s
  \(\delta_0\)), and it lets interior agents in read priority learn to
  keep the slot free while L0 is due or about to be. The share is exact
  in expectation when every L0 file covers the key (Lemma
  D.8\textquotesingle s uniform-key L0 files).
\item
  \textbf{Space:} level \(i\) is charged its \emph{shadowed garbage}
  \(g_i = (1 - \hat{\tilde\rho}_i)B_i\) at the rate \(c_sg_i\,q/\bar q\)
  (§1): the bytes its next compactions would drop at the recent net
  pass-through \(\hat{\tilde\rho}_i\) (trivial moves drop nothing). This
  is an estimate, and it counts only garbage in level \(i+1\) shadowed
  by level \(i\). A version two or more levels below its next-newer
  version is dropped only after that version descends, and no \(g_k\)
  counts it, so \(\sum_ig_i\) undercounts resident garbage whenever
  upper-level keys also live deeper than the next level. Live bytes are
  not attributed: they are the same under every policy (Lemma D.15).
\end{itemize}

\textbf{Proposition D.16 (decomposition).} The per-level charges, plus
the shared hit-read bucket, sum to the global cost rate \(c(t)\) exactly
for the write and read terms. For space, \(\sum_i g_i\) counts only
adjacent shadowing and so undercounts total garbage; it is reported
against measured garbage at run end (OBJ-3).

\emph{Proof.} Every written byte has exactly one source (a flush or one
compaction). Every probe, seek and false-positive block read happens at
exactly one level and is charged in full either there or, under slot
blocking, split between L0 and one other level in shares that sum to 1.
Every hit\textquotesingle s block read goes to the bucket.
\(\blacksquare\)

\subsubsection{§5 Room by mode: predictions recorded before any
run}\label{5-room-by-mode-predictions-recorded-before-any-run}

\begin{itemize}
\tightlist
\item
  \textbf{Read priority.} On \texttt{Assoc} at the old comparators there
  was almost no room in steady state (audit §2). The read levers are the
  L0 trigger --- static at fixed prices, dynamic when prices move
  (Corollary D.12) --- and depth, which is static (Corollary A.7); under
  skew the interior profile is a third read lever, also static
  (Proposition D.13(ii)). Two further levers act during a run, both
  through the single compaction slot: the L0 agent compacting early
  while the slot is idle and reads are heavy (Pathway A §4 (e)), and
  interior levels keeping the slot free while L0 is due or about to be
  (Pathway A §4 (f), priced by the slot-blocking charge of §4). Their
  size is unmeasured. Room is expected mainly on phased or changing
  workloads.
\item
  \textbf{Write priority.} A static profile gains \(O((1-\eta)^2)\)
  (Theorem A.2(iii), where \(\eta\) is a ratio of merged volumes, not
  merge survival), plus what timing merges against garbage can add,
  bounded by Proposition B.1″ and its native-relative form (Pathway B
  §3). On \texttt{Assoc} garbage is not scarce in the measured phase,
  since every Put overwrites, but native compaction already drops
  77--92\% of it, much of it high in the tree (merge survival
  \(X/(S+O)\) at L1 0.84--0.93, D-4). First overwrites meet their loaded
  versions in the deepest levels, where a drop saves the fewest later
  stages. The room is what holding a level adds to high drops beyond
  native; Gate N3 measures it.
\item
  \textbf{Space priority.} On \texttt{Assoc}, \(S\) is 1.03--1.09, so
  garbage is \((S-1)/S\) = 3--8\% of held bytes; even a policy that held
  no garbage at all could remove no more than that. High-garbage
  workloads are needed.
\end{itemize}

\subsubsection{§6 Acceptance}\label{6-acceptance}

{\def\LTcaptype{none} % do not increment counter
\begin{longtable}[]{@{}
  >{\raggedright\arraybackslash}p{(\linewidth - 6\tabcolsep) * \real{0.0292}}
  >{\raggedright\arraybackslash}p{(\linewidth - 6\tabcolsep) * \real{0.2508}}
  >{\raggedright\arraybackslash}p{(\linewidth - 6\tabcolsep) * \real{0.3500}}
  >{\raggedright\arraybackslash}p{(\linewidth - 6\tabcolsep) * \real{0.3500}}@{}}
\toprule\noalign{}
\begin{minipage}[b]{\linewidth}\raggedright
\#
\end{minipage} & \begin{minipage}[b]{\linewidth}\raggedright
Criterion
\end{minipage} & \begin{minipage}[b]{\linewidth}\raggedright
Threshold
\end{minipage} & \begin{minipage}[b]{\linewidth}\raggedright
Instrument
\end{minipage} \\
\midrule\noalign{}
\endhead
\bottomrule\noalign{}
\endlastfoot
OBJ-1 & Flow identity (Lemma D.1, Proposition D.16) & on every arm that
runs the plugin: the attribution log\textquotesingle s per-level write
bytes (flushes to L0, each compaction to its start level, by completion
time), summed over levels and over every interval from \(n_w\) to the
end of the drain --- intervals excluded from replay included --- equal
the evaluator\textquotesingle s flush and compaction bytes for the same
window, with the log opening a partial interval for every level at
\(n_w\) and closing one at the end of the drain; any residual is listed
by job and must consist only of jobs whose completion the two sources
place on different sides of a boundary stamp. The log\textquotesingle s
per-level read counts, taken before the slot-blocking rule moves any of
them and with each level\textquotesingle s hit block reads logged apart
from its false-positive ones, equal the per-level
counters\textquotesingle{} totals over the same window (their difference
between the \(n_w\) stamp and the end of the drain); after the rule,
per-level read charges plus the hit-read bucket still sum to the global
read cost. Space is not checked here: the per-level charge is a garbage
estimate (OBJ-3) & attribution log, event log \\
OBJ-2 & Prices calibrated & all prices in money: device times for
\(c_w, c_f, c_{blk}, c_{sk}\) measured on the node and converted at the
instance price, the storage price \(c_s > 0\), both money prices fixed
in advance (§0.6), and \(\bar q\) per workload, in the fingerprint;
re-measured on any hardware change; every result also reported at
\(c_s/2\) and \(2c_s\) & calibration script \\
OBJ-3 & Space attribution & \(\sum_ig_i\) and measured garbage at run
end reported per arm. \(\sum_ig_i\) targets only adjacent shadowing and
undercounts (§4), so the criterion is on differences: across
\(\Theta_s\) at Gate N2, the ratio\textquotesingle s spread and the
slope of \(\Delta\sum_ig_i\) on \(\Delta\)(measured garbage) are
reported, and a tolerance on the ratio is fixed from them before Gate
N4; at Gate N4 the learner\textquotesingle s paired \(\Delta\sum_ig_i\)
against its comparator must agree in sign with \(\Delta\)(measured
garbage) whenever that difference\textquotesingle s paired interval
excludes zero, and its ratio must lie within that tolerance & reference
compaction, event log \\
OBJ-4 & Per-level read counters & per-level counters taken at the sites
of §4 (the version\textquotesingle s level: \texttt{FilePicker} hits in
\texttt{Version::Get}; seeks per L0 file and per level iterator), so
that a trivially moved file\textquotesingle s reads are charged to its
new level; block reads split into the hit\textquotesingle s read and
false-positive reads (§4); their sums match the global counters within
1\% (the read half of Proposition D.16) & per-level counters, tickers \\
OBJ-5 & Mode recorded & \(\beta\), mode and \(\bar q\) in every
arm\textquotesingle s manifest and fingerprint &
\texttt{metadata.env} \\
OBJ-6 & Unpriced time reported & measured-phase throughput, stall
seconds and controller CPU per arm, paired against the comparator; not
part of \(J_\beta\), but the first two decide the stall rule (Global
acceptance) & db\_bench, event log, per-thread CPU clocks (plugin
thread, trainer process) \\
\end{longtable}
}

\begin{center}\rule{0.5\linewidth}{0.5pt}\end{center}

\subsection{Pathway A --- Actuation: level target multipliers and the L0
trigger}\label{pathway-a--actuation-level-target-multipliers-and-the-l0-trigger}

\subsubsection{§0 Purpose}\label{0-purpose}

Give the controller three moves per level --- compact now, defer, and
defer while expanding the level\textquotesingle s capacity --- through
RocksDB\textquotesingle s own compaction scores, so that RocksDB keeps
sole authority over which files are compacted (history §3.3). The
2026-09-11 motivation, avoiding a depth cascade, is withdrawn (§0.1).

\subsubsection{§1 Mechanism}\label{1-mechanism}

\textbf{New option \texttt{level\_target\_multipliers}} (mutable
column-family option):

\begin{itemize}
\tightlist
\item
  a \texttt{vector\textless{}double\textgreater{}} with one entry per
  level; entry 0 must be 1.0;
\item
  rejected when
  \texttt{level\_compaction\_dynamic\_level\_bytes\ =\ true} or when the
  compaction style is not leveled;
\item
  changeable at run time through \texttt{SetOptions};
\item
  in \texttt{ComputeCompactionScore}, level \(i \ge 1\) is scored as
  \[s_i = \frac{\text{bytes\_not\_compacting}_i}{\texttt{MaxBytesForLevel}(i)\times m_i}.\]
\end{itemize}

\textbf{L0} keeps its native scoring. The controller changes only the
mutable \texttt{level0\_file\_num\_compaction\_trigger} (\(K_0\)),
through \texttt{SetOptions}.

\subsubsection{§2 Actions}\label{2-actions}

For level \(i \ge 1\) the effective multiplier is
\(m_i = \bar m_i\,d_i\): a persistent \textbf{anchor} \(\bar m_i\) (the
level\textquotesingle s capacity) and a short-lived \textbf{timing
factor} \(d_i\) that relaxes back to 1 with time constant
\(\kappa_d\tau_i\).

{\def\LTcaptype{none} % do not increment counter
\begin{longtable}[]{@{}
  >{\raggedright\arraybackslash}p{(\linewidth - 4\tabcolsep) * \real{0.2594}}
  >{\raggedright\arraybackslash}p{(\linewidth - 4\tabcolsep) * \real{0.4324}}
  >{\raggedright\arraybackslash}p{(\linewidth - 4\tabcolsep) * \real{0.2882}}@{}}
\toprule\noalign{}
\begin{minipage}[b]{\linewidth}\raggedright
Action
\end{minipage} & \begin{minipage}[b]{\linewidth}\raggedright
Effect
\end{minipage} & \begin{minipage}[b]{\linewidth}\raggedright
Allowed when
\end{minipage} \\
\midrule\noalign{}
\endhead
\bottomrule\noalign{}
\endlastfoot
hold & nothing: RocksDB proceeds natively under the current \(m_i\) &
always \\
compact & \(d_i \leftarrow \varphi_i/(\bar m_i(1+\epsilon))\), so
\(s_i > 1\) now & \(s_i < 1\) and \(\varphi_i \ge \varphi_{\min}\) \\
defer & \(d_i \leftarrow \varphi_i(1+\epsilon)/\bar m_i\), so
\(s_i < 1\) now & \(s_i \ge 1 - \epsilon\) \\
expand (defer and increase capacity) &
\(\bar m_i \leftarrow \min(\alpha\bar m_i, m_{\max})\),
\(d_i \leftarrow 1\) & \(\bar m_i < m_{\max}\) \\
\end{longtable}
}

\begin{itemize}
\tightlist
\item
  \textbf{Hold is the default.} It means "let RocksDB act natively", so
  a controller that always holds is native RocksDB (ACT-4).
\item
  \textbf{Anchors decay.} When not re-expanded,
  \(\bar m_i \leftarrow 1 + (\bar m_i - 1)e^{-\Delta t/(\kappa_a\tau_i)}\),
  so expansion cannot ratchet (ACT-5).
\item
  \textbf{Bounds.} Every \(m_i\) stays in \([m_{\min}, m_{\max}]\)
  (A-Impl-7).
\item
  \textbf{Masks.} An action not allowed in a state is excluded both when
  the action is chosen and in the learning target (Proposition H.2).
\end{itemize}

\textbf{L0 analogues}, acting on \(K_0\): hold; compact,
\(K_0 \leftarrow \max(2, k_0)\) so L0 is due now; defer,
\(K_0 \leftarrow \min(K_{\text{cap}}, k_0 + 1)\) temporarily; expand,
raise the L0 anchor \(\bar K_0\) by one within \([2, K_{\text{cap}}]\),
decaying back to the configured value.

\textbf{Lemma A.5 (what an action does).} Under A3, "compact" makes
level \(i\) eligible at the next score computation and "defer" makes it
ineligible. Neither guarantees admission: RocksDB admits the eligible
level with the highest score when a compaction slot is free.

\emph{Proof.} A3 and RocksDB\textquotesingle s ranking of eligible
levels by score. \(\blacksquare\)

\emph{Consequence.} Actions at different levels interact through score
ranking. That is why neighbour and global state are inputs (H §2), and
why each level is charged for what actually ran (Proposition D.16), not
for what it requested.

\subsubsection{§3 Theory}\label{3-theory}

\textbf{Theorem A.1 (no release while held; corrected).} If
\(\varphi_i(t) < m_i(t)\) throughout \([t_0, t_1]\), level \(i\) is not
eligible and releases nothing in that interval. When it is released at
fill \(\varphi_i \ge m_i\), RocksDB compacts it until its score falls
below 1, so it sends down about \((\varphi_i - m_i)C_i\) source bytes
(to within one file), of which about \(\tilde\rho_i(\varphi_i -
m_i)C_i\) net bytes land below (trivially moved files land whole; here
\(\tilde\rho_i\) is the pass-through of the released files, which after
a hold can differ from the window mean in either direction: a backlog
clears partly by trivial moves, but in D-12 deferrals during the
post-load backlog turned would-be moves into merges, audit 2026-09-23,
§1).

\emph{Proof.} A3, and RocksDB re-scores after each compaction and keeps
picking the level while its score is at least 1 and highest.
\(\blacksquare\)

\emph{Correction.} The 2026-09-11 statement charged a burst of
\(\kappa C_i\). At \(m_i = 1\) the burst is \((\kappa-1)C_i\) (audit
§3).

\emph{Containment (fluid, A5).} A burst of \(b\) net bytes landing at
level \(i+1\) makes level \(j > i\) due only if \(b\) exceeds the
combined headroom \(\sum_{k=i+1}^{j}(m_kC_k - B_k)^+\). Populated depth
grows only if a burst passes the last populated level\textquotesingle s
headroom \((m_LC_L - B_L)^+\). At T=10 on \texttt{Assoc} the last
level\textquotesingle s nominal target (about 16 GiB) is far above its
content (about 1 GB), so no burst can deepen the tree there.

\textbf{Theorem A.2 (fanout conservation; corrected and extended).}

\begin{itemize}
\tightlist
\item
  (i) \(\prod_{i=0}^{L-1}f_i = B_L/(K_0F)\), independent of
  \(m_1, \dots, m_{L-1}\) (Corollary D.9\textquotesingle s product,
  which uses only the definitions of the \(f_i\) and none of that
  corollary\textquotesingle s assumptions).
\item
  (ii) With weights \(v_i = a_i - t_i\), the merged bytes leaving level
  \(i\) per user byte (Lemma D.7), held fixed, and a common overlap
  constant \(c\), the fanout part of the write cost, \(c\sum_iv_if_i\),
  is minimised under (i) at \(f_i \propto 1/v_i\), with minimum
  \(cL\,(B_L/(K_0F))^{1/L}(\prod_i
  v_i)^{1/L}\).
\item
  (iii) For \(v_i = \eta^i\), that is a constant ratio
  \(\eta = v_{i+1}/v_i =
  \tilde\rho_i(1-\xi_{i+1})/(1-\xi_i)\) of merged volumes (not the merge
  survival \(\eta_i = X/(S+O)\) of Pathway B §2), the relative gain over
  equal fanouts is
  \[1 - \frac{L\,\eta^{(L-1)/2}(1-\eta)}{1-\eta^{L}} = O\big((1-\eta)^2\big).\]
\end{itemize}

{\def\LTcaptype{none} % do not increment counter
\begin{longtable}[]{@{}
  >{\raggedleft\arraybackslash}p{(\linewidth - 6\tabcolsep) * \real{0.0933}}
  >{\raggedleft\arraybackslash}p{(\linewidth - 6\tabcolsep) * \real{0.0233}}
  >{\raggedleft\arraybackslash}p{(\linewidth - 6\tabcolsep) * \real{0.5367}}
  >{\raggedleft\arraybackslash}p{(\linewidth - 6\tabcolsep) * \real{0.3267}}@{}}
\toprule\noalign{}
\begin{minipage}[b]{\linewidth}\raggedleft
\(\eta\)
\end{minipage} & \begin{minipage}[b]{\linewidth}\raggedleft
\(L\)
\end{minipage} & \begin{minipage}[b]{\linewidth}\raggedleft
optimum / equal fanouts
\end{minipage} & \begin{minipage}[b]{\linewidth}\raggedleft
available gain
\end{minipage} \\
\midrule\noalign{}
\endhead
\bottomrule\noalign{}
\endlastfoot
0.95 & 4 & 0.9984 & 0.2\% \\
0.90 & 4 & 0.9931 & 0.7\% \\
0.80 & 5 & 0.9519 & 4.8\% \\
0.60 & 5 & 0.7807 & 21.9\% \\
\end{longtable}
}

\emph{Proof.} (i) Corollary D.9. (ii) With a multiplier on the product
constraint, stationarity gives \(v_if_i = \lambda\) for every \(i\). The
constraint then gives \(\lambda = (B_L/(K_0F)\cdot\prod_iv_i)^{1/L}\),
and the minimum of \(\sum v_if_i\) is \(L\lambda\). (iii) With
\(v_i = \eta^i\), \(\prod_i v_i = \eta^{L(L-1)/2}\), and equal fanouts
give \((B_L/(K_0F))^{1/L}\sum_i\eta^i\); divide. The first-order terms
in \(1-\eta\) cancel. \(\blacksquare\)

\emph{Assumption behind (ii) and (iii).} The weights are held fixed
while the profile varies, but they depend on it. Levels \(\ge 1\) hold
one version per key, so an overwrite is absorbed when the newer version
arrives: a larger level \(i\) holds more keys, the merges into it
(sourced at \(i-1\)) drop more, and \(\rho_{i-1}\) falls, and with it
the weight \(v_i\) of level \(i\)\textquotesingle s own merges; fewer
duplicates are left for the merges out of level \(i\), so \(\rho_i\) may
rise. The trivial-move shares can move too (in one simulation with a
round-robin picker, by more than \(\rho\): enlarging level \(i\) lowered
\(\xi_{i-1}\) and raised \(\xi_i\); under min-overlap picking, the
pinned \texttt{kMinOverlappingRatio}, they barely moved). The overlap
constant moves with the fanouts. (ii) is therefore the optimum for fixed
weights. The joint optimum needs survival measured as a function of the
profile. \(\Theta_s\) supplies it only at its own profiles: three of
them are uniform scalings, which move \(m_i\) and \(m_{i+1}\) together,
so their separate effects on \(\rho_i\) (a larger level \(i+1\) lowers
it) cannot be told apart, and its survival-weighted profile is computed
from weights measured at uniform 1. Treat that profile as one step of a
fixed-point iteration: re-measure the weights under it and recompute; if
the profile moves by more than its measurement interval, add the second
profile to \(\Theta_s\) before the comparator is fixed (§0.6).

\emph{Corrections.} The optimal profile needs some multipliers below 1,
which the 2026-09-11 scale (\(s \ge 1\)) could not express (audit §3);
multipliers below 1 are now allowed. The profile is static, so its gain
belongs to the comparator (Corollary C.3). The accounting in (ii) is
Lemma D.7\textquotesingle s, not a convention.

\textbf{Corollary A.3 (2026-09-11: level removal is never profitable).}
\emph{Superseded.} Its premise, that removing a level needs \(s = T\),
failed at T=10 where \(s = 2\) removed one (audit §3). Whether fewer
levels pay depends on the priority (Proposition D.13(iv)), and removal
is static in any case (Corollary A.7).

\textbf{Corollary A.4 (2026-09-11: write-optimal size ratio).}
\emph{Superseded} by Proposition D.13(iv), in which the optimum depends
on the measured overlap constant and the priority.

\textbf{Lemma A.6 (depth never falls during a run).} Under A4 and
without deletes, the index of the deepest non-empty level never
decreases.

\emph{Proof.} Compactions move bytes from level \(i\) to \(i+1\) or
within a level. The deepest non-empty level can only receive bytes or
push them deeper, and without deletes no compaction removes all of its
data. \(\blacksquare\)

\textbf{Corollary A.7 (level removal is a static effect).} A controller
can prevent depth growth but cannot remove a level during a run.
Removing a level --- the largest read benefit static expansion showed
(audit §3; at \(s = 1.5\), with depth unchanged, reads also fell 6.3\%
at T=2 and 3.4\% at T=10, the static hit shift of Proposition D.13(ii))
--- is a property of the configuration the tree was loaded under (base
size, multipliers at load time), and so belongs to the static
comparator.

\subsubsection{§4 What actuation can and cannot
deliver}\label{4-what-actuation-can-and-cannot-deliver}

\textbf{Static, so the comparator has it too:} the equal-fanout or
survival-weighted profile (Theorem A.2), depth fixed at load (Corollary
A.7), \(K_0\) at fixed prices (Proposition D.11), and the interior
profile\textquotesingle s read effect (Proposition D.13(ii)).

\textbf{Dynamic, so this is the controller\textquotesingle s only case.}
Each item is a hypothesis tested without a learner at Gate N3:

\begin{itemize}
\tightlist
\item
  (a) tracking \(K_0^\star\) and the profile as prices change with the
  workload (Corollary D.12);
\item
  (b) timing a release against the neighbours: overlap per source byte
  is about \(T\varphi_{i+1}/\varphi_i\) (Lemma D.8), and a burst may be
  arriving from above;
\item
  (c) holding a level so its merges drop more garbage, bounded by the
  garbage native leaves and by how much higher the drops can move
  (Pathway B §3);
\item
  (d) preventing depth growth (Theorem A.1);
\item
  (e) L0 compacting early (the L0 "compact" action) while the compaction
  slot is idle and reads are heavy, and holding while it is busy;
\item
  (f) interior levels holding their releases while L0 is due or about to
  be, so that L0 is not kept waiting for the slot (the slot-blocking
  charge, Pathway D §4).
\end{itemize}

(e) and (f) are read levers that act during a run. A fixed trigger or
profile can express neither, because each depends on the
slot\textquotesingle s state at the moment of the decision.

\subsubsection{§5 Implementation}\label{5-implementation}

\begin{itemize}
\tightlist
\item
  \textbf{A-Impl-1. Scale inside the score only.} Never route the
  multiplier through \texttt{max\_bytes\_for\_level\_base}. That option
  is also L0\textquotesingle s byte threshold and L1\textquotesingle s
  base, so routing through it would silently change L0\textquotesingle s
  trigger. The implementation folds \(m_i\) into
  \texttt{MaxBytesForLevel(i)} through the fork\textquotesingle s
  existing per-level \texttt{capacity\_scales\_} (which already leaves
  level 0 unscaled). That gives the score of §1 exactly, and makes
  A-Impl-3 automatic. The only other reader of \texttt{MaxBytesForLevel}
  that affects file picking is the round-robin input expansion in
  \texttt{compaction\_picker\_level.cc}, which the pinned
  \texttt{kMinOverlappingRatio} never runs. Tested (ACT-1):
  L0\textquotesingle s score is unchanged under any multiplier vector.
\item
  \textbf{A-Impl-2. Static ladder.}
  \texttt{level\_compaction\_dynamic\_level\_bytes\ =\ false} is pinned,
  and option validation rejects multipliers otherwise (A1).
\item
  \textbf{A-Impl-3. Pending-compaction estimate.}
  \texttt{EstimateCompactionBytesNeeded} computes pending bytes from
  level targets. It must apply \(m_i\) too, or RocksDB holds two
  different notions of what is due, and the pending-bytes slowdown and
  speed-up logic reads unscaled targets. It reads
  \texttt{MaxBytesForLevel}, so the implementation of A-Impl-1 satisfies
  this (verified at \texttt{25468bbaa}). Re-check the soft and hard
  pending-bytes limits.
\item
  \textbf{A-Impl-4. When a change takes effect.} Verified at
  \texttt{25468bbaa} (\texttt{DBImpl::SetOptions}): the call appends a
  new Version without a MANIFEST write, which recomputes every score,
  before it returns, waiting its turn behind any flush or compaction
  install. An action is therefore in the published scores at once, with
  or without writes.
\item
  \textbf{A-Impl-5. Cost of \texttt{SetOptions}.} Verified at
  \texttt{25468bbaa}: each call appends that Version and installs a new
  SuperVersion under the DB mutex, then writes and renames a full
  OPTIONS file with the mutex released (\texttt{WriteOptionsFile};
  RocksDB has no switch to skip it). The OPTIONS file is not an SST, so
  it enters none of \(W\), \(H\) or \(J_\beta\) (D §1); its time cost is
  measured by ACT-2 and reported, not priced. Call \texttt{SetOptions}
  only when a value changes (hold makes no call), send all
  levels\textquotesingle{} changes in one call, and cap the call rate.
\item
  \textbf{A-Impl-6. L0 trigger range.} \([2, K_{\text{cap}}]\) with
  \(K_{\text{cap}} =
  \min(\lfloor C_1/F\rfloor, K_{\text{slow}} - 1)\) (A3′). Gate 1 found
  triggers 8 and 16 identical at an 8 MiB base.
\item
  \textbf{A-Impl-7. Multiplier bounds.}
  \(m_i \in [m_{\min}, m_{\max}]\), fixed in advance; \(m_{\min} = 0.5\)
  is suggested. \(m_{\max} \le 2.0\) unless larger values are measured,
  since 2.0 is only the edge of the tested static grid (audit §3). Level
  targets must also never shrink going down the tree:
  \(m_{i+1}C_{i+1} \ge m_iC_i\), i.e. \(m_{i+1}T \ge m_i\). RocksDB
  assumes this ordering (\texttt{version\_set.cc} asserts it over
  non-empty levels), and the fork\textquotesingle s capacity code
  already rejects vectors that break it. At \(T = 2\) the bounds
  \([0.5, 2]\) alone do not guarantee it, so option validation rejects
  such a vector, and the controller\textquotesingle s masks exclude any
  action that would produce one.
\item
  \textbf{A-Impl-8. Attribution and fallback.} Log every change with
  level, old and new value, action, decision id and reason. If the
  controller is absent or its weights are invalid, reset \(m \equiv 1\)
  and \(K_0\) to its configured value, and log it as a fallback.
\item
  \textbf{A-Impl-9. Ruled out:
  \texttt{max\_bytes\_for\_level\_multiplier\_additional}.} It is an
  integer vector, its effect accumulates down the tree, and it is
  ignored under dynamic level sizing (as in 2026-09-11).
\item
  \textbf{A-Impl-10. Native parity.} With \(m \equiv 1\) and \(K_0\) at
  its configured value, the patched binary must reproduce stock RocksDB
  (ACT-4).
\end{itemize}

\subsubsection{§6 Acceptance}\label{6-acceptance-1}

{\def\LTcaptype{none} % do not increment counter
\begin{longtable}[]{@{}
  >{\raggedright\arraybackslash}p{(\linewidth - 6\tabcolsep) * \real{0.0340}}
  >{\raggedright\arraybackslash}p{(\linewidth - 6\tabcolsep) * \real{0.1293}}
  >{\raggedright\arraybackslash}p{(\linewidth - 6\tabcolsep) * \real{0.4083}}
  >{\raggedright\arraybackslash}p{(\linewidth - 6\tabcolsep) * \real{0.4083}}@{}}
\toprule\noalign{}
\begin{minipage}[b]{\linewidth}\raggedright
\#
\end{minipage} & \begin{minipage}[b]{\linewidth}\raggedright
Criterion
\end{minipage} & \begin{minipage}[b]{\linewidth}\raggedright
Threshold
\end{minipage} & \begin{minipage}[b]{\linewidth}\raggedright
Instrument
\end{minipage} \\
\midrule\noalign{}
\endhead
\bottomrule\noalign{}
\endlastfoot
ACT-1 & Option correctness & the fork\textquotesingle s test file
\texttt{level\_target\_multipliers\_test} is green (Debug build), and
the same checks pass on the Release binary in the preflight: L0 score
invariance, rejection under dynamic sizing and of vectors whose targets
shrink going down, scaled pending estimate, recompute after
\texttt{SetOptions} & fork gtest; preflight check script over
\texttt{db\_bench} output and LOG \\
ACT-2 & \texttt{SetOptions} overhead & foreground throughput cost ≤ 1\%
at the planned call rate, paired, OPTIONS-file writes included; a
diagnostic of overhead, not part of \(J_\beta\) & db\_bench \\
ACT-3 & Action fidelity & a requested \(m_i\) appears in the published
score within one control interval for ≥ 99\% of changes & policy log,
score observer \\
ACT-4 & Native parity & patched binary at \(m \equiv 1\) inside the
oracle-parity envelope of stock RocksDB (the checks of the 2026-08-22
gate) & paired evaluator \\
ACT-5 & No ratchet & anchors return to within \(\epsilon\) of 1 after
\(5\kappa_a\tau_i\) without re-expansion & policy log \\
\end{longtable}
}

The 2026-09-11 criteria A-0 to A-5 are retired (history §18.1).

\begin{center}\rule{0.5\linewidth}{0.5pt}\end{center}

\subsection{Pathway G --- Propagation across levels (turnover
normalisation)}\label{pathway-g--propagation-across-levels-turnover-normalisation}

\subsubsection{§0 Purpose and scope}\label{0-purpose-and-scope}

Deep levels complete few outcomes per run, so a learner at a deep level
has too little data. RusKey reports the same for its FLSM-tree: training
every level separately failed from Level 3 onward {[}RusKey, §7{]}. This
pathway lets deep levels use what upper levels learn, without copying
decisions downward and without stopping deep levels from learning.

\textbf{Scope.} Only \emph{interior} levels, \(1 \le i \le L-1\), which
have a level below them, can be pooled. L0 is tiered and has its own
trigger action (A3′). The last level has no output level, and its only
role is whether the tree deepens (Lemma A.6). Both are separate,
unpooled agents (Pathway H).

\subsubsection{§1 Related work: RusKey\textquotesingle s policy
propagation}\label{1-related-work-ruskeys-policy-propagation}

What RusKey does {[}RusKey, §5{]}:

\begin{itemize}
\tightlist
\item
  \textbf{Why deep levels lack data.} Compaction at a deeper level is
  exponentially less frequent, because FLSM merges a whole level at a
  time.
\item
  \textbf{Case 1, uniform bits-per-key.} Only Level 1\textquotesingle s
  policy \(K_1\) is learned, with the other levels held fixed; every
  deeper level is then set to \(K_1\).
\item
  \textbf{Case 2, Monkey filters.} Levels 1 and 2 are learned until
  stable. Deeper levels are then set by a closed-form recurrence (their
  Lemma 5.1), derived by minimising a per-level cost in which
  false-positive rates grow by \(T\) per level, and rounded to an
  integer.
\item
  \textbf{After the transfer,} deep levels are not learned.
\item
  \textbf{Also in Lerp:} one independent DDPG agent per level; actions
  limited to \(K-1\), \(K\), \(K+1\); reward
  \(\alpha\cdot(\text{level latency}) + (1-\alpha)\cdot
  (\text{end-to-end latency})\).
\end{itemize}

Two passages in the paper are inconsistent with the rest; describe the
method from the lemma and the figure, not from these sentences:

\begin{itemize}
\tightlist
\item
  The Case 2 bullet says Level \(i\) has \emph{lower} bits-per-key than
  Level \(i+1\), while the Monkey description and the lemma have
  false-positive rates \emph{rising} with depth.
\item
  The Fig. 9 text says the policy gets lazier with depth, while the
  figure (10, 8, 3, 1) and the lemma\textquotesingle s worked example
  (9, 7, 3, 1) show \(K\) falling with depth, i.e. more aggressive
  deeper.
\end{itemize}

\textbf{How this pathway differs.}

{\def\LTcaptype{none} % do not increment counter
\begin{longtable}[]{@{}
  >{\raggedright\arraybackslash}p{(\linewidth - 4\tabcolsep) * \real{0.1907}}
  >{\raggedright\arraybackslash}p{(\linewidth - 4\tabcolsep) * \real{0.3946}}
  >{\raggedright\arraybackslash}p{(\linewidth - 4\tabcolsep) * \real{0.3946}}@{}}
\toprule\noalign{}
\begin{minipage}[b]{\linewidth}\raggedright
\end{minipage} & \begin{minipage}[b]{\linewidth}\raggedright
RusKey (Lerp)
\end{minipage} & \begin{minipage}[b]{\linewidth}\raggedright
This pathway
\end{minipage} \\
\midrule\noalign{}
\endhead
\bottomrule\noalign{}
\endlastfoot
What propagates & the chosen \(K\), copied (Case 1) or extrapolated by a
closed form (Case 2) & a value estimate in level-free units; each level
still decides from its own state \\
Justification & the closed-form optimum of a per-level analytic cost &
each level faces the same problem on its own clock (Proposition G.1) at
its own prices (Proposition G.2); differences between levels are
measured inputs \\
Deep levels after propagation & fixed & keep learning; their own
correction grows with their own data (Lemma G.5) \\
Timing & sequential: learn L1 (and L2) to convergence, then transfer
once & concurrent and continuous \\
Why deep data is scarce & deeper compactions are exponentially rarer
(whole-level merges) & with RocksDB\textquotesingle s file-by-file
compaction, job counts are similar across levels; completed
\emph{turnovers} are about \(T/\tilde\rho\) times rarer per level \\
Knob & runs per level, \(K \in [1, T]\), in FLSM & target multiplier on
leveled RocksDB (\(K = 1\)); L0 trigger separately \\
Neighbours & level statistics and all levels\textquotesingle{} policies
in the state & the burst about to land from above, as an explicit
forecast input \\
Reward & \(\alpha\)-mix of level and end-to-end latency & attributed
priced cost plus a counterfactual neighbour charge (H §3) \\
\end{longtable}
}

\textbf{Shared ideas, cited to RusKey:} one agent per level, and actions
limited to one step at a time.

\subsubsection{§2 Two propositions: the level clock and the price per
turnover}\label{2-two-propositions-the-level-clock-and-the-price-per-turnover}

Definitions are in §1.1: pass-through \(\rho_i\) over merges,
trivial-move share \(\xi_i\), net pass-through \(\tilde\rho_i\), overlap
\(o_i\), inflow \(\lambda_i\), outflow \(\mu_i\), turnover time
\(\tau_i\), level clock \(\theta_i\), survival to level \(i\) \(\pi_i\).
In addition, per level: \(e^f_i\) filter probes per Get at level \(i\),
\(e^b_i\) false-positive block reads per Get at level \(i\) (the
hit\textquotesingle s block read goes to the hit-read bucket, Pathway D
§4), \(\nu_i\) runs seeked per scan at level \(i\) (1 if the level has a
file at or after the seek key), and the holding price
\(\sigma_i = c_sN_i/(c_w\bar q)\), with \(N_i\) the operations served
per turnover of level \(i\) (§1.1): the cost of holding one byte while
level \(i\) turns over once (space is charged per operation served,
Pathway D §1), in units of the cost of writing one byte.

\textbf{Proposition G.1 (each level runs on its own clock).}

\begin{itemize}
\tightlist
\item
  (a)
  \(\dfrac{d\varphi_i}{d\theta_i} = \dfrac{\lambda_i - \mu_i}{\bar\lambda_i}\).
\item
  (b) In steady state, \(\tau_{i+1} = (T/\tilde\rho_i)\,\tau_i\).
\end{itemize}

\emph{Proof.} (a) \(dB_i/dt = \lambda_i - \mu_i\); substitute
\(B_i = \varphi_iC_i\) and \(dt = \tau_i\,d\theta_i\) with
\(\tau_i = C_i/\bar\lambda_i\). (b) In steady state
\(\bar\mu_i = \bar\lambda_i\), and by the definition of \(\tilde\rho_i\)
as the net share of level \(i\)\textquotesingle s outflow that lands in
level \(i+1\) (merged bytes that survive plus trivially moved bytes),
\(\bar\lambda_{i+1} = \tilde\rho_i\bar\mu_i\). Hence
\(\tau_{i+1}/\tau_i = (C_{i+1}/C_i)(\bar\lambda_i/\bar\lambda_{i+1}) =
T/\tilde\rho_i\). \(\blacksquare\)

In plain terms: on its own clock every interior level fills at the same
average rate, one level\textquotesingle s worth per turnover. What
differs between levels is how many seconds one tick takes --- about
\(T\) times longer per level down.

\emph{Evidence that jobs are not the scarce quantity.} In D-7 at T=2,
L1, L2 and L3 released 935, 1,079 and 927 times (history §14.20). That
arm compacted early, so treat the numbers as indicative. At the old
measured-phase ingest rate (about 7 MB/s), T=10 gives \(\tau \approx 2\)
s at L1, 23 s at L2 and 230 s at L3, so the old 160 s runs completed no
turnover at L3.

\textbf{Proposition G.2 (price per turnover).} In steady state, level
\(i\)\textquotesingle s attributed part of \(C_\beta\) (Pathway D §4:
its write and read shares, and for space its shadowed-garbage charge,
which is not a share of \(C_\beta\)\textquotesingle s space term)
accumulated over one turnover, divided by \(c_wC_i\) (the price of
writing one level\textquotesingle s worth of bytes), is
\[\hat c_i = \beta_W(1-\xi_i)(\rho_i + o_i) + \frac{\beta_R\big(\tilde R^f e^f_i + \tilde R^b e^b_i + \tilde R_{sk}\nu_i + \tilde y_i\big)}{\pi_i} + \beta_S\,\sigma_i\,(1-\tilde\rho_i)\bar\varphi_i,\]
with the workload price ratios, the same at every level,
\[\tilde R^f = \frac{q_{pt}c_f}{c_w\bar\lambda_1},\qquad \tilde R^b = \frac{q_{pt}c_{blk}}{c_w\bar\lambda_1},\qquad \tilde R_{sk} = \frac{q_{sc}c_{sk}}{c_w\bar\lambda_1},\]
and \(\tilde y_i = y_i/(c_w\bar\lambda_1)\), where \(y_i\) is the mean
unweighted L0 read cost per second charged to level \(i\) by the
slot-blocking rule (Pathway D §4).

\emph{Proof.} Over one turnover the level releases
\(\bar\mu_i\tau_i = C_i\) source bytes in steady state. A share
\(1-\xi_i\) of them is merged and writes \((1-\xi_i)C_i(\rho_i + o_i)\)
bytes; trivially moved bytes write nothing (Lemma D.7). Reads charged to
the level cost \(\beta_R\tau_i[q_{pt}(c_fe^f_i +
c_{blk}e^b_i) + q_{sc}c_{sk}\nu_i + y_i]\). The level\textquotesingle s
attributed garbage \((1-\tilde\rho_i)B_i\) (Pathway D §4) costs
\(\beta_S(c_s/\bar q)
(1-\tilde\rho_i)\bar\varphi_iC_iN_i\) over the \(N_i\) operations of the
turnover. Divide by \(c_wC_i\) and use \(\tau_i/C_i = 1/\bar\lambda_i =
1/(\bar\lambda_1\pi_i)\). \(\blacksquare\)

In plain terms: the write part has the same form at every level, with
\(\rho_i\), \(\xi_i\) and \(o_i\) as measured inputs. The levels differ
through those inputs, through the read share they carry (\(e_i\), and
\(\pi_i\)) and through the holding price \(\sigma_i\), which grows about
\(T\) times per level down. These become \emph{inputs} to a shared model
rather than things it must learn.

\subsubsection{§3 The propagation law}\label{3-the-propagation-law}

\textbf{Definition (law).} For an interior level \(j\) in the pool
\(\mathcal P\), with normalised state \(\hat x_j\) and action \(a\), the
action value (expected future normalised reward, rewards being negative
costs) is
\[Q_j(\hat x_j, a) = b(\hat x_j, a) + f_\theta(\hat x_j, a) + \delta_j(\hat x_j, a),\]
where \(b\) is the analytic prior (code, H §7), \(f_\theta\) is a
learned correction shared by the pool, and \(\delta_j\) is a small
per-level correction. Both learned parts start at zero (cold start, H
§4). The law has five rules:

\begin{itemize}
\tightlist
\item
  \textbf{(G-i) Pooled training in level-free units.} Every pooled
  level\textquotesingle s transitions enter one training set after
  conversion: reward
  \(\hat r = -(\text{level cost over the transition})/(c_wC_j)\), the
  level cost being H §3\textquotesingle s (attributed cost plus
  neighbour charge), and discount
  \(\gamma_j = \exp(-\Delta N/(n_HN_j))\), where \(\Delta N\) is the
  number of operations served over the transition, \(N_j\) the
  operations per turnover of level \(j\) (§1.1), and \(n_H\) the
  look-ahead in turnovers, the same at every level. The discount runs on
  operations, like \(J_\beta\): on the wall clock, any action that
  delays operations --- a stall, or compaction that slows the foreground
  --- would lower every later discounted cost by about the
  delay\textquotesingle s share of the horizon, a gain \(J_\beta\) does
  not contain. \(f_\theta\) is fitted on this pooled set. Pooling
  converted transitions is the "shared experience" option; the
  level-specific inputs keep it from imposing one
  level\textquotesingle s garbage rate on another. Every length "in
  turnovers" in Pathways G and H --- a transition\textquotesingle s
  length, the time since a release, the exploration spend,
  \(n_{\text{turn}}\) --- is counted in operations, \(\Delta N/N_j\); it
  equals \(\Delta t/\tau_j\) while the operation rate is steady, which
  is the sense of the level clock \(\theta_i\) in Proposition G.1.
\item
  \textbf{(G-ii) Per-level correction.} \(\delta_j\) is fitted only on
  level \(j\)\textquotesingle s own transitions, pulled toward zero with
  a strength of \(n_0\) turnovers (Lemma G.5). Each transition is
  counted by its duration in turnovers, because decisions inside one
  turnover are strongly correlated.
\item
  \textbf{(G-iii) Own decisions.} Each level acts on its own \(Q_j\) and
  its own state. Nothing is copied between levels.
\item
  \textbf{(G-iv) Cadence.} Level \(j\) decides every \(N_j/k\)
  operations, with \(k\) fixed in advance, so every transition serves
  the same number of operations and carries the same discount,
  \(e^{-1/(n_Hk)}\). During a write stop no decision falls due;
  RocksDB\textquotesingle s own triggers act. Decisions that fall due
  together go into one \texttt{SetOptions} call (A-Impl-5). Amended
  2026-10-02: L0, which is not pooled, decides once per flush, every
  \(N_0/K_0^{\text{cfg}}\) operations, since its state changes only at
  flushes and L0 compactions; its transitions carry
  \(e^{-1/(n_H K_0^{\text{cfg}})}\). The \texttt{SetOptions} rate cap
  (A-Impl-5) must stay below the shortest control interval in wall time,
  or ACT-3 cannot hold: at \(k = 10\), L0 decided every 30--64 ms on the
  2026-10-02 rules smoke, against a 100 ms cap.
\item
  \textbf{(G-v) Membership.} A level joins \(\mathcal P\) only if it
  passes the admission test (§4). A level that fails keeps its own
  model.
\end{itemize}

\textbf{Normalised state} of level \(j\), all dimensionless:

\begin{itemize}
\tightlist
\item
  own: \(\varphi_j\), \(\bar m_j\), \(d_j\), \(s_j\), and time since its
  last release in turnovers;
\item
  neighbours: \(\varphi_{j\pm1}\), \(m_{j\pm1}\);
\item
  incoming burst, for \(j \ge 2\):
  \(b_j = \tilde\rho_{j-1}(\varphi_{j-1} - m_{j-1})^+/T\), the share of
  level \(j\)\textquotesingle s capacity about to land from above;
\item
  slot contention \(\chi_j = (\hat\omega_j, \zeta)\): over the last
  decision interval, level \(j\)\textquotesingle s mean wait per release
  for the compaction slot, counted in operations served meanwhile and in
  units of the decision interval, and the share of the
  interval\textquotesingle s operations served while the slot was busy
  (Proposition G.4). A job\textquotesingle s wait runs from the later of
  the level\textquotesingle s score reaching 1 and its previous job
  ending, to the job starting; \(\hat\omega_j\) is carried forward from
  the last interval with a release. These are averages over the
  interval, not snapshots: the slot changes state several times a
  second, far faster than any interior level\textquotesingle s decision
  interval;
\item
  inflow ratio \(\ell_j = \lambda_j^{\text{recent}}/\bar\lambda_j\):
  whether inflow is running above or below its average;
\item
  queue position: the number of due levels (L0 included) whose score
  exceeds \(\max(s_j, 1+\epsilon)\), i.e. how many levels
  RocksDB\textquotesingle s score order would try before a compaction at
  \(j\), now or after "compact" (Lemma A.5: an action changes
  eligibility, not admission); and where the running
  compaction\textquotesingle s start level lies relative to \(j\) (none,
  above, same, below);
\item
  backlog: RocksDB\textquotesingle s pending-compaction estimate, with
  the multipliers applied (A-Impl-3), over the held SST bytes \(H\); and
  L0\textquotesingle s distance to the slowdown trigger,
  \((K_{\text{slow}} - k_0)/K_{\text{slow}}\). Both are the same at
  every level. The estimate is divided by \(H\), not by the soft
  pending-bytes limit, which here is 64 GiB against about 3 GiB held and
  would read near 0;
\item
  two levels down, and the bottom: \(\varphi_{j+2}/m_{j+2}\) (the last
  level\textquotesingle s value when \(j+2 = L\); 0 with an absent flag
  when \(j+2 > L\)), and the last level\textquotesingle s free room
  \((m_LC_L - B_L)^+/(m_LC_L)\). A burst that overflows level \(j+1\)
  reaches \(j+2\) next (Theorem A.1, containment), and the last
  level\textquotesingle s room decides whether the tree deepens (Lemma
  A.6);
\item
  level terms: \(\hat\rho_j\), \(\hat\xi_j\), the overlap constant
  \(c_j\), \(\pi_j\), \(e^f_j\), \(e^b_j\), \(\nu_j\), \(\sigma_j\);
\item
  workload price ratios \(\tilde R^f\), \(\tilde R^b\),
  \(\tilde R_{sk}\);
\item
  the priority vector \(\beta\), and L0\textquotesingle s \(k_0/K_0\).
\end{itemize}

\textbf{Proposition G.3 (normalisation does not change a
level\textquotesingle s decisions).} Define level
\(j\)\textquotesingle s objective with the discount of (G-i). Then its
action values in currency equal \(c_wC_j\) times its normalised action
values, so every ranking of actions, and every greedy choice, is the
same in both units.

\emph{Proof.} Every reward of level \(j\) is divided by the same
positive constant, and the discount depends on \(\Delta N/N_j\) in both.
The Bellman equation is therefore scaled by that constant, and so is its
unique fixed point. Ranking is invariant under positive scaling.
\(\blacksquare\)

\emph{The modelling choice this rests on.} Discounting per turnover,
counted in operations, rather than per second is a choice, justified by
Proposition G.1: a level\textquotesingle s outcomes arrive on its own
clock. It gives deep levels a longer horizon in seconds, which the
2026-09-11 controller lacked (0.98 per second, about 50 s of look-ahead
at every level).

\textbf{Proposition G.4 (when two levels can share one value function).}
Write level \(j\)\textquotesingle s normalised state as
\(\hat x_j = (z_j, \vartheta_j)\): \(z_j\) the dynamic part (fills,
multipliers, time since release, neighbours, burst, inflow ratio,
contention, queue position, backlog, the fills two levels down and at
the bottom) and \(\vartheta_j\) the level terms and prices, which are
fixed within a run and include the overlap constant \(c_j\). Hold the
other agents\textquotesingle{} policies fixed. Let levels \(i\) and
\(j\) satisfy:

\begin{itemize}
\tightlist
\item
  (c1) conditional on \((\hat x, a)\), the law of the next normalised
  state and of the transition\textquotesingle s length in turnovers is
  the same function at both levels. In particular \(\hat x\) is Markov:
  the next state depends on the rest of the tree only through
  \(\hat x\);
\item
  (c2) the reward over a transition is the same function of
  \((\hat x, a, \hat
  x')\) at both levels. The attributed cost is, by the per-transition
  form of Proposition G.2\textquotesingle s accounting, once
  \(c_j \in \vartheta_j\). The neighbour charges of H §3 are the same
  function only if both levels\textquotesingle{} neighbours are pool
  members; at the pool\textquotesingle s edges (next to L1 or the last
  level) the difference enters \(\varepsilon_r\) below;
\item
  (c3) one compaction job moves or delivers a negligible share of the
  level (source file and incoming batch \(\ll C_i\)).
\end{itemize}

Then the two levels\textquotesingle{} optimal normalised action values
are the same function of \(\hat x\).

\emph{Proof.} Under (c1) and (c2), and with the same admissible actions
as a function of \(\hat x\) (the masks\textquotesingle{}
\(\varphi_{\min}\), \(\epsilon\), \(m_{\min}\), \(m_{\max}\) and
\(\alpha\) equal at both levels), the two normalised problems are one
semi-Markov decision problem. With the cadence of (G-iv) every
transition serves \(N_j/k\) operations, so its discount is
\(\gamma = e^{-1/(n_Hk)} < 1\). The Bellman operator is a contraction
with a unique fixed point. (c3) is what makes (c1) plausible on the
level\textquotesingle s own clock (A5). \(\blacksquare\)

\emph{When the conditions hold only approximately.} Let the expected
rewards given \((\hat x, a)\) differ by at most \(\varepsilon_r\), and
let \(|r| \le R\). If the laws of the next state differ by at most
\(\varepsilon_P\) in total variation, uniformly in \((\hat x, a)\), then
\(\lVert Q^\star_i - Q^\star_j\rVert_\infty \le
\varepsilon_r/(1-\gamma) + 2\varepsilon_PR/(1-\gamma)^2\) {[}Kearns \&
Singh, simulation lemma{]}, with
\(1/(1-\gamma) = n_Hk + \tfrac12 + O(1/(n_Hk))\). Total variation has no
scale: a variable averaged over the level\textquotesingle s own decision
interval (\(\zeta\), \(\hat\omega_j\), \(\ell_j\)) has a law that
tightens with depth, which drives \(\varepsilon_P\) toward 1. The useful
form assumes \(\gamma V^\star_j\) is \(L\)-Lipschitz in the next state
and uses the Wasserstein-1 distance \(\varepsilon_W\) instead:
\(\lVert Q^\star_i - Q^\star_j\rVert_\infty \le
(\varepsilon_r + L\varepsilon_W)/(1-\gamma)\). Lipschitz continuity is
itself an assumption; the mask thresholds can make \(V^\star\) jump.

\emph{Contention is where (c1) is only approximate.} With
\texttt{max\_background\_jobs\ =\ 2} and
\texttt{max\_background\_flushes\ =\ max\_background\_compactions\ =\ −1}
(db\_bench\textquotesingle s defaults; the pipeline sets only the first,
fingerprint field \texttt{bg2}), \texttt{DBImpl::GetBGJobLimits} gives
one flush slot, \(\max(1, \lfloor 2/4\rfloor)\), and one compaction
slot, \(\max(1, 2-1)\). The speed-up path can only set the compaction
count to one, so it is one in every state. The bottom-priority pool
counts against the same limit and db\_bench starts it with no threads,
and \texttt{max\_subcompactions} is 1
(\texttt{db/db\_impl/db\_impl\_compaction\_flush.cc},
\texttt{GetBGJobLimits} and \texttt{MaybeScheduleFlushOrCompaction};
unchanged from the pinned base \texttt{7ea2d73} to the recorded commit
\texttt{25468bbaa}). A level that becomes due while the slot is busy, or
while a level with a higher score is due, waits. The wait cannot be made
level-free: the slot changes state on the wall clock, several times a
second at every depth, so on level \(j\)\textquotesingle s clock its
dynamics run about \(T/\tilde\rho\) times faster per level down. What
makes pooling possible is that the wait is short on a deep
level\textquotesingle s clock. Let \(\omega_j\) be the mean of
\(\hat\omega_j\): the mean wait per release in units of level
\(j\)\textquotesingle s decision interval \(N_j/k\). The contention part
of \(\varepsilon_W\) (per unit of \(L\)) between two levels is then at
most about \(\omega_i + \omega_j\) (a heuristic: a delay moves outflow
across a decision boundary with probability about its length over the
interval, if decisions are not synchronised with releases). \(\omega_j\)
falls roughly \(T/\tilde\rho\)-fold per level down and is largest at L1,
one reason L1 is expected to fail. It is measured (Gate N0, item 3) and
bounded in the admission test (§4).

The content is in (c1) and (c2). The admission test screens them on
native runs; PROP-1b re-checks (c1) under the controller.

\textbf{Lemma G.5 (how a level\textquotesingle s trust in its own data
grows).} Let \(\delta_j\) be linear in \(k\) standardised, mutually
orthogonal features, fitted by least squares with penalty
\(n_0\lVert w\rVert^2\) on \(n_j\) own samples. Each fitted coefficient
equals the own-data least-squares coefficient times \(n_j/(n_j + n_0)\).

\emph{Proof.} With \(X^\top X = n_jI\), the penalised solution is
\(w = (n_jI +
n_0I)^{-1}X^\top y = \frac{n_j}{n_j+n_0}\cdot\frac{X^\top y}{n_j}\), and
\(X^\top y/n_j\) is the unpenalised solution. \(\blacksquare\)

In plain terms: with no data of its own a level uses the shared model;
after \(n_0\) turnovers of its own, half of its correction comes from
its own data. Keep \(\delta_j\) to a few features so it can be fitted
from little data.

\subsubsection{§4 Admission test (the collapse
test)}\label{4-admission-test-the-collapse-test}

From static runs (\(m \equiv 1\), one configuration per (workload,
\(T\), \(K_0\))), compute for each interior level, on the
level\textquotesingle s own clock:

\begin{itemize}
\tightlist
\item
  the fill sampled at fixed points of the level clock (every \(1/k\)
  turnover), and the fill at release;
\item
  bytes released per turnover, in units of \(C_i\);
\item
  \((1-\xi_i)(\rho_i + o_i)\) and \(\tilde\rho_i\) per turnover;
\item
  the normalised inflow ratio \(\ell_i\);
\item
  the mean slot wait per release, \(\omega_i\), in decision intervals
  (G.4).
\end{itemize}

\((1-\xi_i)(\rho_i + o_i)\) and \(\tilde\rho_i\) are model inputs;
comparing them is a support screen, so that \(f_\theta\) is not
extrapolated beyond the levels it was fitted on.

Statistics whose scale is set by the file size rather than the level are
(c3) diagnostics, reported but not compared: jobs per turnover, and time
between jobs in turnovers, both fall by about \(T/\tilde\rho\) per level
whatever the controller does.

A level joins the pool only if it is shown equivalent to the reference
level, named in advance (the shallowest candidate interior level). For
each compared statistic \(s\), the differences in mean and in the 10th
and 90th percentiles, in the statistic\textquotesingle s own units, must
have 90\% intervals inside \(\pm\delta_{\text{adm},s}\). Margins are
fixed in advance in those units; for statistics in fill units they are
no smaller than the file granularity \(F_{\text{sst}}/C_i\) of the
shallower level. The upper 95\% bound on \(\omega_i\) must be below
\(\omega_{\max}\). Intervals come from a moving-block bootstrap over
whole turnovers (all events of a turnover together), resampled within
each run, with block length \(b\) fixed in advance. Requiring every
statistic to pass is an intersection--union test, so the chance of
admitting a level that differs by more than its margin on any one of
them is at most 5\% without a multiplicity correction (asymptotically;
the intervals are bootstrap intervals). The minimum number of turnovers
per level, \(n_{\min}\), is fixed in advance by simulation: resample the
reference level\textquotesingle s own turnovers into two pseudo-levels,
which must pass with probability at least 0.8; shift one statistic at a
time to its margin, the others unshifted, and in each case the pair must
pass with probability at most 0.05.

A Kolmogorov--Smirnov distance is reported but is not the criterion. It
has no scale, so it rejects tight distributions that differ by less than
one file. Its plug-in value is also biased upward: at 30 turnovers, the
upper 95\% bootstrap bound between identical distributions has a median
of 0.43, so identical levels would almost never pass.

A test that merely fails to reject a difference is not evidence of
sameness: with few turnovers it passes for lack of data. Static runs
measure native dynamics, while (c1) must hold under the controller, so
membership is re-checked on the learner\textquotesingle s own
transitions by the conditional check PROP-1b.

Expected failures: L1, which is fed in L0-sized batches and whose job
size is comparable to \(C_1\) at small base sizes; the last level; and
the level above the last whenever the last level is far from its target,
since its fanout \(f_{L-1} = B_L/(m_{L-1}C_{L-1})\) is set by the last
level\textquotesingle s fill (at T=10 on \texttt{Assoc},
\(f_3 \approx 0.6\) against \(f_2 = 10\)), so the T=10 pool may be L2
alone. Gate N1\textquotesingle s pilot native runs (PREREGISTRATION
D-16; the 2026-09-11 programme\textquotesingle s artifacts are out of
scope) decide every level that completes at least \(n_{\min}\) turnovers
in them. Levels with fewer stay undecided and unpooled (D-19, which
withdrew their later decision on Gate N2\textquotesingle s native arms).
On the 2026-10-01 pilots no cell had a pool (D-19).

\textbf{Scope decision (2026-09-29): propagation is claimed only where
the tree is deep enough to pool.} A pool needs at least two admitted
levels. L1 and the last level never qualify, and the level above the
last qualifies only when the last level is near its target, so a pool
needs the last populated level at L5 or deeper, or at L4 when L4 is near
its target. With about 3 GiB of \texttt{Assoc} data that means T=2
(populated to L8, candidates L2--L6); T=6 is decided by the admission
test. At T=10 the tree reaches L4 with L4 far below target, so the pool
is at most L2 alone and nothing propagates. Every level then keeps its
own model, and L3, which completes about one turnover per run, learns
little and stays close to its prior and to hold. Pooling a level that
failed admission (starting it from \(f_\theta\)) is not done: at T=10 it
would amount to copying L2\textquotesingle s policy into a level whose
merges write about 1.6 bytes per byte moved against L2\textquotesingle s
11, and it would use \(f_\theta\) outside the levels it was fitted on.
More data at T=10 (about 17 GiB, so that L4 nears its target and L3 can
join L2) is the escalation path, to be decided by one static run at that
size before any other.

\subsubsection{§5 What propagation does and does not
deliver}\label{5-what-propagation-does-and-does-not-deliver}

By Lemma D.10, an interior level\textquotesingle s own filter probes do
not depend on its multiplier (except through the L0 run count: directly
for L1, weakly for every level through the shared compaction slot), and
a scan seeks each level with a file at or after the seek key once
however full it is. Raising \(m_i\) moves hits up from deeper levels
(Proposition D.13(ii)): level \(i\) takes over hit reads, which go to
the hit-read bucket and are charged to no level (Pathway D §4), and the
probes at levels \(i+1..L\) fall. The saving, which can be material
under skew, lands in other levels\textquotesingle{} attributed costs;
the neighbour charges of H §3 do not carry it, since their one-step
prediction models bursts and overlap only. It is a static effect of the
profile and belongs to \(\Theta_s\). Pooled interior agents therefore
mainly save write and space cost --- timing releases against neighbour
fills and incoming bursts, and holding data so merges drop more garbage
--- and, in read priority, keep the compaction slot free while L0 is due
or about to be (Pathway A §4 (f)), which the slot-blocking charge
prices.

The dynamic read levers are the L0 trigger, L0\textquotesingle s use of
an idle slot, interior levels\textquotesingle{} yielding of the slot,
and depth. Only slot yielding involves the pool. The interior
profile\textquotesingle s read effect is static. In read priority any
read gain comes mainly from the L0 agent, from slot timing and from not
deepening the tree. This agrees with the 2026-09-11 record.

\subsubsection{§6 Acceptance}\label{6-acceptance-2}

{\def\LTcaptype{none} % do not increment counter
\begin{longtable}[]{@{}
  >{\raggedright\arraybackslash}p{(\linewidth - 6\tabcolsep) * \real{0.0635}}
  >{\raggedright\arraybackslash}p{(\linewidth - 6\tabcolsep) * \real{0.2269}}
  >{\raggedright\arraybackslash}p{(\linewidth - 6\tabcolsep) * \real{0.5444}}
  >{\raggedright\arraybackslash}p{(\linewidth - 6\tabcolsep) * \real{0.1452}}@{}}
\toprule\noalign{}
\begin{minipage}[b]{\linewidth}\raggedright
\#
\end{minipage} & \begin{minipage}[b]{\linewidth}\raggedright
Criterion
\end{minipage} & \begin{minipage}[b]{\linewidth}\raggedright
Threshold
\end{minipage} & \begin{minipage}[b]{\linewidth}\raggedright
Instrument
\end{minipage} \\
\midrule\noalign{}
\endhead
\bottomrule\noalign{}
\endlastfoot
PROP-1 & Admission recorded & collapse test run and pool membership
recorded per (workload, \(T\)) before any learner run & Gate N1 \\
PROP-1b & Membership under control & for each pooled level, a one-step
model of the normalised transition (\(\hat x' \mid \hat x, a\): fill
change, bytes released, time to next release) fitted on the pool has, on
that level\textquotesingle s held-out learner transitions, a mean
residual whose 90\% block-bootstrap interval lies inside
\(\pm\delta_{\text{kern}}\); and adding a level indicator lowers
held-out error by less than \(\delta_{\text{kern}}\). A level that fails
leaves the pool at the next weight push, and the removal is logged. The
check is conditional, not marginal, because each level\textquotesingle s
own policy changes its marginals & training log \\
PROP-2 & Transfer helps prediction & on each pooled
level\textquotesingle s held-out transitions, reported per level, the
pooled model\textquotesingle s normalised prediction error is below that
of a model trained only on that level\textquotesingle s data, paired
over seeds & training log \\
PROP-3 & No harm & \(J_\beta\) with the pooled deep-level agents acting
is not worse than with those levels held (paired interval upper bound ≤
0, or a margin fixed in advance); their attributed cost is reported
alongside & paired evaluator \\
PROP-4 & Ablation & propagation on versus off (off = per-level models on
own data), per mode & Gate N6 \\
PROP-5 & Related work & the §1 comparison with RusKey appears in the
paper\textquotesingle s related work & paper \\
\end{longtable}
}

\begin{center}\rule{0.5\linewidth}{0.5pt}\end{center}

\subsection{Pathway H --- RL architecture (Programme
1)}\label{pathway-h--rl-architecture-programme-1}

\subsubsection{§0 Overview}\label{0-overview}

\begin{verbatim}
 workload (db_bench mixgraph / suite, Pathway B)
        |
        v
 RocksDB (pinned base + level_target_multipliers)  <-- SetOptions, batched (A-Impl-5)
        |  state read in-process                              ^
        v                                                      |
 C++ controller plugin (separately fingerprinted library)  ----+
   - L0 agent (K0), interior agents (pooled, Pathway G), last-level agent
   - masks, analytic prior b, fallback to native
        |  transitions out                ^  weights in, every Δ_push
        v                                 |
 Python trainer: shared model f_θ + per-level corrections δ_j, Double DQN
\end{verbatim}

Inference happens where the state is. Training happens off the critical
path. Every learned quantity starts at zero at the start of each run
(cold start).

\subsubsection{§1 Agents}\label{1-agents}

\begin{itemize}
\tightlist
\item
  \textbf{L0 agent.} Actions on \(K_0\) (Pathway A §2). Not pooled
  (A3′). It holds the main read lever (Proposition D.11).
\item
  \textbf{Interior agents, levels \(1..L-1\).} Actions on \(m_i\).
  Levels that pass the admission test share one model (Pathway G); the
  others keep their own.
\item
  \textbf{Last-level agent.} Actions hold and expand only; compact and
  defer have no meaning for a level far below its target. Its job is to
  stop the tree deepening when the last level nears its target (Theorem
  A.1). On \texttt{Assoc} at 10M and T=10 it is idle, since the last
  level is far below its target (Theorem A.1, containment).
\item
  \textbf{All agents} see the same priority vector \(\beta\) and
  workload prices, and all start from hold.
\end{itemize}

\subsubsection{§2 State}\label{2-state}

\begin{itemize}
\tightlist
\item
  \textbf{Interior agents:} the normalised state of G §3.
\item
  \textbf{L0 agent:} \(k_0/K_0\), \(\bar K_0\), flush rate over its
  mean, L1 fill \(\varphi_1\), the Get-to-write and scan-to-write rate
  ratios, the price ratios \(\tilde R\), \(\beta\), and the compaction
  slot\textquotesingle s busy share over the last interval (L0 is the
  level most exposed to the single slot, G.4); and, as for interior
  agents (G §3), queue position, backlog and L2\textquotesingle s fill
  \(\varphi_2/m_2\). For lever (e) of Pathway A §4 it also sees whether
  the slot is idle now (the queue position\textquotesingle s "none"),
  and the active memtable\textquotesingle s fill over
  \texttt{write\_buffer\_size} (how soon the next flush adds an L0 file;
  RocksDB property \texttt{rocksdb.cur-size-active-mem-table}).
\item
  \textbf{Last-level agent:} \(\varphi_L/m_L\), headroom
  \((m_LC_L - B_L)/C_L\), the incoming burst, queue position and
  backlog.
\end{itemize}

\textbf{Removed} from the 2026-09-11 state: the eleven stall-era
pressure inputs, the Lagrange multipliers, and the scan-work and
space-estimate inputs (the scan metric sat at its floor; the space
estimate was retired by D-3).

\subsubsection{§3 Reward: attributed cost plus a counterfactual
neighbour
charge}\label{3-reward-attributed-cost-plus-a-counterfactual-neighbour-charge}

For agent \(i\) over a decision interval,
\[r_i = -\frac{c^\beta_i(\Delta t) + X_{i+1} + X_{i-1}}{c_wC_i},\]

\begin{itemize}
\tightlist
\item
  \(c^\beta_i(\Delta t)\) is level \(i\)\textquotesingle s attributed,
  priority-weighted cost over the interval (Pathway D §4);
\item
  \(X_{i\pm1} = \bar V_{i\pm1}(x'_{i\pm1}\mid a_i) - \bar V_{i\pm1}(x'_{i\pm1}\mid
  \text{hold})\) is the change in the neighbour\textquotesingle s
  expected future cost caused by level \(i\) taking \(a_i\) instead of
  holding. It is evaluated with the neighbour\textquotesingle s current
  value estimate, converted to currency by \(c_wC_{i\pm1}\), and a
  one-step prediction of the neighbour\textquotesingle s state: the
  burst \(i\) releases or keeps (Theorem A.1) and the overlap that
  implies (Lemma D.8).
\item
  \(X = 0\) when \(a_i\) is hold.
\item
  The divisor is a constant per agent. For the L0 agent it is
  \(C_0 := K_0^{\text{cfg}}F\), the bytes L0 holds at its configured
  trigger, with \(F\) as measured at \(n_w\) and then held fixed; for
  the last-level agent it is \(C_L\). A constant divisor leaves the
  agent\textquotesingle s ranking of actions unchanged (Proposition
  G.3). The live \(K_0F\) would not, since \(K_0\) is the L0
  agent\textquotesingle s own action.
\end{itemize}

The charges do not model read shifts. A level that expands no longer
pays the block reads of the hits it takes over from deeper levels, since
those go to the hit-read bucket (D §4), but it is credited nothing for
the filter probes saved below it (G §5). The read effect of the profile
is static and belongs to \(\Theta_s\). The one read effect an interior
level has during a run, keeping L0 waiting for the compaction slot, is
in its attributed cost \(c^\beta_i\) through the slot-blocking charge (D
§4).

In plain terms: a level pays for what it spends, and also for what its
choice makes its neighbours spend later. That is how a compaction that
is good for one level but bad for the next one is discouraged.

\textbf{Proposition H.1 (why a counterfactual charge, not a shaping
term).}

\begin{itemize}
\tightlist
\item
  (a) Adding \(\gamma\Phi(s') - \Phi(s)\) to the reward, for any
  function \(\Phi\) of the state, cannot change which policy is optimal
  {[}Ng et al.{]}. So charging a level for how its
  neighbour\textquotesingle s value changes over time --- a function of
  state that level \(i\) observes --- cannot make it take the neighbour
  into account.
\item
  (b) The difference reward
  \(D_i = G(a_i, a_{-i}) - G(\text{hold}, a_{-i})\), with \(G\) the
  global cost and \(a_{-i}\) the other agents\textquotesingle{} actions,
  changes by exactly \(G\)\textquotesingle s change when agent \(i\)
  alone changes its action:
  \(D_i(a') - D_i(a) = G(a', a_{-i}) - G(a, a_{-i})\) {[}Wolpert \&
  Tumer{]}.
\end{itemize}

\emph{Proof.} (a) Ng et al., Theorem 1; for transitions of variable
duration use Proposition H.4. (b) The subtracted term does not depend on
\(a_i\). \(\blacksquare\)

\emph{The approximation.} The reward above approximates \(D_i\) in two
ways: it restricts \(G\)\textquotesingle s change to level \(i\) and its
two neighbours, and it estimates the neighbours\textquotesingle{} change
with learned values. ARCH-3 measures the one-step state prediction the
charge relies on. The design follows counterfactual credit assignment in
cooperative multi-agent learning {[}Wolpert \& Tumer; Foerster et
al.{]}. It is not RusKey\textquotesingle s \(\alpha\)-mix of level and
end-to-end latency.

Two further gaps separate these rewards from \(J_\beta\) even if every
value were exact. First, \(\bar V_{i\pm1}\) is the
neighbour\textquotesingle s value under its own reward, which includes
the neighbour\textquotesingle s charge for its effect on level \(i\); so
\(X_{i\pm1}\) carries back part of level \(i\)\textquotesingle s own
future cost, which level \(i\) also pays directly. A value head fitted
to the neighbour\textquotesingle s attributed cost alone avoids this
echo; which of the two is used is fixed in advance (§0.6), and a Gate N6
ablation compares the two. Second, a difference reward aligns one
agent\textquotesingle s unilateral change with \(G\) (Proposition
H.1(b)), but agents that each maximise their own return --- discounted
per turnover of their own level (G-i), while \(J_\beta\) is an
undiscounted sum over the measured phase --- can settle where no agent
can improve its own return and \(J_\beta\) is still not minimal. The
per-level returns are therefore a training surrogate; every claim is
scored on \(J_\beta\) (Global acceptance).

\subsubsection{§4 Learning rule}\label{4-learning-rule}

\begin{itemize}
\tightlist
\item
  \textbf{Double DQN} {[}van Hasselt et al.; Mnih et al.{]} on
  normalised rewards with the turnover discount of (G-i); Huber loss, a
  target network and gradient clipping.
\item
  \textbf{Masked target} (fixes the defect found in the audit):
  \(y = r + \gamma\,Q_{\text{target}}\big(s', \arg\max_{a' \in \mathcal A(s')}
  Q_{\text{online}}(s', a')\big)\), the \(\arg\max\) taken only over the
  actions allowed in \(s'\).
\item
  \textbf{Residual form.} \(Q = b + f_\theta + \delta_j\) (Pathway G).
  \(f_\theta\) and \(\delta_j\) start at zero; \(b\) is code.
\item
  \textbf{Replay} keeps only transitions with valid attribution;
  fallback intervals are excluded.
\end{itemize}

\textbf{Proposition H.2 (an unmasked target overestimates).} Let
\(\mathcal A(s')
\subsetneq \mathcal A\) for some reachable \(s'\). If the target
maximises over all of \(\mathcal A\), its fixed point \(\tilde Q\)
satisfies \(\tilde Q \ge Q^\star\) pointwise, where \(Q^\star\) is the
masked fixed point. The inequality is strict at every \((s, a)\) that
reaches, with positive probability, a state \(s'\) in which a forbidden
action has the highest value.

\emph{Proof.} A maximum over a superset is at least the maximum over the
subset, so the unmasked Bellman operator dominates the masked one
pointwise. Both are monotone contractions, so iterating from the same
start keeps the order, and the limits satisfy \(\tilde Q \ge Q^\star\).
Strictness follows from the positive-probability state. \(\blacksquare\)

In plain terms: the learner credits states with the value of actions it
will never be allowed to take there.

\subsubsection{§5 Timing, exploration, and the measured
phase}\label{5-timing-exploration-and-the-measured-phase}

\begin{itemize}
\tightlist
\item
  \textbf{Settle, then measure} (adopted from the parallel stream;
  amended 2026-09-30, \texttt{docs/PREREGISTRATION.md} D-13). After the
  load, \texttt{db\_bench} issues no operation until
  RocksDB\textquotesingle s \texttt{WaitForCompact}, flushes included,
  returns. It then holds for \(h_w\) = 10 s, during which every
  level\textquotesingle s score must stay below 1 and \(k_0 < K_0\). The
  measured phase starts at \(n_w\), the first \texttt{mixgraph}
  operation, on the tree the arm\textquotesingle s own load left. Every
  arm, native included, is scored from \(n_w\) (A8), and the controller
  starts at \(n_w\), so no controller action falls outside the scored
  phase. An arm not settled at the end of the hold is invalid. The
  earlier rule drained the backlog under live traffic and set \(n_w\)
  from pilot runs with a margin, leaving the operations before \(n_w\)
  unscored; this one needs no pilots and scores every \texttt{mixgraph}
  operation.
\item
  \textbf{Default is hold.} Exploration departs from hold with
  probability \(p_x\) per decision and never takes a masked action.
  \(p_x\) keeps a floor after training, so the policy cannot freeze. The
  exploration spend is counted per level in turnovers.
\item
  \textbf{Units.} The prior\textquotesingle s action differences are
  expressed in the same normalised units as \(Q\). This fixes the D-7
  mismatch, where prior terms of about ±2 were added to \(Q\) values in
  the thousands.
\end{itemize}

\textbf{Proposition H.3 (exploration has a cost; the 2026-09-11
Proposition B.3).} Let the best static configuration \(\theta^\star\) be
constant over the run, and suppose the optimal policy for the control
problem is realisable as a constant action. Then any online controller
\(\pi\) that does not know \(\theta^\star\) and explores
non-degenerately has, over a finite horizon,
\(J(\pi) = J(\theta^\star) + \mathcal R\) with \(\mathcal R > 0\).

\emph{Proof.} As 2026-09-11: under the hypothesis,
\(\pi\)\textquotesingle s cost is \(\theta^\star\)\textquotesingle s
cost plus regret on every interval where its action differs, and any
schedule with a positive probability of a suboptimal action on a set of
positive measure contributes strictly positive regret. \(\blacksquare\)

The hypothesis --- that a constant action is optimal --- is exactly what
Conjecture B.3′ doubts.

\emph{Consequence.} Exploration is reported as a cost, and the first
\(N_x\) turnovers after \(n_w\) are logged separately.

\textbf{Proposition H.4 (shaping with variable durations; the 2026-09-11
Proposition D.1).} For semi-Markov transitions of duration \(\tau\), the
term \(F = \gamma^{\tau}\Phi(s') - \Phi(s)\) leaves optimal policies
unchanged; the single-step form \(\gamma\Phi(s') - \Phi(s)\) does not
when \(\tau\) varies.

\emph{Proof.} As 2026-09-11, following Ng et al. with the discount of
each transition raised to its duration. \(\blacksquare\)

It applies to any shaping term added later.

\subsubsection{§6 Where inference and training
run}\label{6-where-inference-and-training-run}

\begin{itemize}
\tightlist
\item
  \textbf{C++ plugin inside the RocksDB process.} It reads the state
  in-process at decision time, evaluates the per-level models, and
  applies \texttt{SetOptions} in one batched call. There is no socket
  round trip per decision, so the stale-answer rejection of D-12
  (27--39\% of answers) cannot occur.
\item
  \textbf{Python trainer.} It receives transitions, trains, and pushes
  weights every \(\Delta_{\text{push}}\) (fixed in advance). The
  deployed policy lags training by at most \(\Delta_{\text{push}}\). The
  learning rule is off-policy, so it tolerates the lag. The weight
  version is logged with every decision.
\item
  \textbf{A separately fingerprinted library.} Changing the controller
  does not change the database binary\textquotesingle s hash, so the
  static comparator stays valid (audit §6). ARCH-5 checks that the
  plugin at \(m \equiv 1\) reproduces native.
\item
  \textbf{Fallback.} Plugin absent or weights invalid: \(m \equiv 1\)
  and \(K_0\) at its configured value (A-Impl-8).
\end{itemize}

\subsubsection{\texorpdfstring{§7 Analytic prior
\(b\)}{§7 Analytic prior b}}\label{7-analytic-prior-b}

\(b(\hat x, a)\) is the expected change in normalised cost over one
turnover from taking \(a\) instead of hold, computed from Lemmas
D.7--D.10 and Proposition G.2 with the current measured inputs. It
prices:

\begin{itemize}
\tightlist
\item
  a compaction now, at the current overlap (Lemma D.8), plus the
  slot-blocking charge it would incur if L0 falls due while its job
  holds the slot (Pathway D §4), from L0\textquotesingle s fill and the
  active memtable\textquotesingle s fill;
\item
  a deferral, by the garbage it keeps (Pathway D §4) and the burst it
  builds (Theorem A.1);
\item
  an expansion, by its space bound (Lemma D.14);
\item
  for the L0 agent, an early compaction, by the extra L0 → L1 overlap it
  writes (Proposition D.11) against the probes and seeks saved while the
  slot is idle.
\end{itemize}

It is clipped to \(\pm b_{\max}\) in normalised units. It is code, not
trained weights, so cold start is preserved.

\subsubsection{§8 Removed from the 2026-09-11
architecture}\label{8-removed-from-the-2026-09-11-architecture}

\begin{itemize}
\tightlist
\item
  Lagrange multipliers, dual ascent and windowed hinges: Programme 1 has
  no hard constraints.
\item
  The guard, which moves to Programme 2. RocksDB\textquotesingle s own
  slowdown and stop triggers and pending-bytes limits remain the only
  safety mechanism.
\item
  Socket-based inference and the snapshot staleness check.
\item
  The per-second discount; the stall-era state inputs.
\end{itemize}

\subsubsection{§9 Acceptance}\label{9-acceptance}

{\def\LTcaptype{none} % do not increment counter
\begin{longtable}[]{@{}
  >{\raggedright\arraybackslash}p{(\linewidth - 6\tabcolsep) * \real{0.0482}}
  >{\raggedright\arraybackslash}p{(\linewidth - 6\tabcolsep) * \real{0.2008}}
  >{\raggedright\arraybackslash}p{(\linewidth - 6\tabcolsep) * \real{0.4820}}
  >{\raggedright\arraybackslash}p{(\linewidth - 6\tabcolsep) * \real{0.2490}}@{}}
\toprule\noalign{}
\begin{minipage}[b]{\linewidth}\raggedright
\#
\end{minipage} & \begin{minipage}[b]{\linewidth}\raggedright
Criterion
\end{minipage} & \begin{minipage}[b]{\linewidth}\raggedright
Threshold
\end{minipage} & \begin{minipage}[b]{\linewidth}\raggedright
Instrument
\end{minipage} \\
\midrule\noalign{}
\endhead
\bottomrule\noalign{}
\endlastfoot
ARCH-1 & Learning health & TD error in normalised units bounded, final
value below twice the reward scale; no NaN; residual magnitudes reported
& training log \\
ARCH-2 & Masks & unit test of the masked target green
(\texttt{rl\_agent/tests/test\_agent.py}); the training
log\textquotesingle s masked-target audit counts zero target argmaxes
over a masked action; zero masked actions executed & tests, training
log, policy log \\
ARCH-3 & Counterfactual input & one-step prediction of neighbour fill
after an action versus realised, error reported per level pair; if above
a tolerance fixed in advance, that pair uses the local reward only &
policy log \\
ARCH-4 & Exploration and no freeze & exploration share and cost reported
per level; the chosen action varies with state (flip rate and
correlation with \(\varphi\) reported where a level has ≥ 10 turnovers)
& policy log \\
ARCH-5 & Plugin parity & plugin at \(m \equiv 1\) inside the
oracle-parity envelope & paired evaluator \\
ARCH-6 & Inference agreement & C++ inference and Python evaluation of
the same weights agree on ≥ 99.9\% of logged decisions & replay \\
\end{longtable}
}

\begin{center}\rule{0.5\linewidth}{0.5pt}\end{center}

\subsection{Pathway B --- Workloads and
garbage}\label{pathway-b--workloads-and-garbage}

\subsubsection{§0 Purpose}\label{0-purpose-1}

Programme 1 is measured across a suite of workloads, for two reasons. No
static setting is best everywhere, which is what gives a controller
something to add across workloads (Definition C.5). And each priority
mode needs a workload on which it has room (Pathway D §5).

\subsubsection{§1 The suite}\label{1-the-suite}

{\def\LTcaptype{none} % do not increment counter
\begin{longtable}[]{@{}
  >{\raggedright\arraybackslash}p{(\linewidth - 4\tabcolsep) * \real{0.3267}}
  >{\raggedright\arraybackslash}p{(\linewidth - 4\tabcolsep) * \real{0.3267}}
  >{\raggedright\arraybackslash}p{(\linewidth - 4\tabcolsep) * \real{0.3267}}@{}}
\toprule\noalign{}
\begin{minipage}[b]{\linewidth}\raggedright
Family
\end{minipage} & \begin{minipage}[b]{\linewidth}\raggedright
Construction
\end{minipage} & \begin{minipage}[b]{\linewidth}\raggedright
Label in the paper
\end{minipage} \\
\midrule\noalign{}
\endhead
\bottomrule\noalign{}
\endlastfoot
UDB \texttt{Assoc} (primary) {[}Cao et al.{]} & \texttt{db\_bench}
mixgraph with the published fit (B1), with the deviations of §2 & "Assoc
key distribution and operation mix at the project\textquotesingle s
record size" \\
Other Cao et al. mixes & published fits where available, e.g. the
ZippyDB mix & realistic \\
Zipfian mixes in the style of YCSB A, B, C and F {[}Cooper et al.{]} &
\texttt{db\_bench} or YCSB & standard \\
Read-heavy / write-heavy / mixed & operation mix swept & synthetic \\
Garbage level & overwrite share and key-space size swept, giving low and
high \(g_{\text{flow}}\) & synthetic \\
Hot ranges & \texttt{keyrange\_num} \(\in \{5, 30, 100\}\) &
synthetic \\
Deletes (B3) & \texttt{mix\_delete\_ratio} \(\in \{0, 3, 6\}\%\) &
synthetic \\
Phases (B2) & two phases first, more later; in-process
\texttt{db\_bench} phase patch & synthetic \\
\end{longtable}
}

Every run records its family and parameters in the fingerprint. A static
comparator measured on one family is never a comparator for a policy
measured on another (the workload form of the knob-parity rule).

\subsubsection{§2 Implementation notes carried
forward}\label{2-implementation-notes-carried-forward}

\begin{enumerate}
\def\labelenumi{\arabic{enumi}.}
\tightlist
\item
  \textbf{(Done 2026-09-20.)} \texttt{WORKLOAD\_SKEW} selects the
  family. Deviations from the published \texttt{Assoc} fit:
  \texttt{value\_theta} is 925.5, not 0, keeping the mean value size at
  the project\textquotesingle s 960 bytes so the level ladder and the
  \(T\) sweep stay comparable; \texttt{mix\_max\_value\_size} is 65536,
  because \texttt{db\_bench} applies it as
  \texttt{val\_size\ \%\ value\_max} and the default 1024 wraps 6.85\%
  of draws down to as little as one byte, pulling the mean to 890.2.
\item
  \textbf{Phases (B2)} as a \texttt{db\_bench} phase patch. Trap:
  \texttt{QueryDecider::Initiate} resets its range total but appends to
  \texttt{type\_} and \texttt{ratio\_} without clearing them, and
  \texttt{GetType} returns the first threshold the draw falls below
  (\texttt{tools/db\_bench\_tool.cc}). A second \texttt{Initiate} for a
  new phase leaves the first phase\textquotesingle s thresholds in
  front, so the old mix keeps running and nothing reports it. Clear both
  vectors, or build a new \texttt{QueryDecider}, at each phase boundary.
  Phase boundaries go in the manifest and fingerprint.
\item
  \textbf{Deletes (B3)} as a fourth operation type in the same patch,
  appended at index 3 after Get, Put and Seek. \texttt{GetType} returns
  a position in the ratio vector, so reordering the vector would remap
  every existing mix; a delete ratio of 0 leaves the existing thresholds
  unchanged.
\item
  \textbf{Per-level survival.} Per compaction record \(S\), \(O\), \(X\)
  and the source level, excluding trivial moves (whose output equals
  input and would bias the ratios toward 1), and record trivially moved
  bytes per source level separately. Derive \(\rho_i\), \(o_i\), \(t_i\)
  and \(\xi_i\), dropped bytes, and the 2026-09-11 merge survival
  \(\eta_i = X/(S+O)\).
\item
  \textbf{Hindsight-oracle runner.} Per phase, run the static sweep and
  record the best static cost; compose them into \(\mathcal G\).
\end{enumerate}

\subsubsection{§3 Theory}\label{3-theory-1}

\textbf{Proposition B.1′ (pooled survival floor; the part of the
2026-09-11 Theorem B.1 that is proved).} Suppose total compaction input
over the run is at least the bytes the user wrote. Then pooled merge
survival --- total compaction output over total compaction input --- is
at least \(1/S_{\text{flow}}\).

\emph{Proof.} Each compaction\textquotesingle s output is its input
minus what it drops. Each obsolete byte is dropped at most once, so
total drops are at most \(N_{\text{written}} -
N_{\text{live}}\). Hence output/input \(\ge 1 - (N_{\text{written}} -
N_{\text{live}})/\text{input} \ge N_{\text{live}}/N_{\text{written}}\),
using input \(\ge N_{\text{written}}\). \(\blacksquare\)

\textbf{Status of the 2026-09-11 ceiling.} Theorem B.1 went on to
evaluate the formula \(\sum_i\eta^i\), whose \(\eta\) is a ratio of
merged volumes as in Theorem A.2(iii), at \(\eta = 1/S_{\text{flow}}\),
and called the result a ceiling on the achievable reduction of
\(W - 1\). The floor justifies neither step. It bounds pooled merge
survival \(X/(S+O)\), whose denominator holds the overlap, so a dropped
byte lowers the volume ratio \(1 + o_i\) times as much as it lowers
merge survival. And it limits pooled survival, not how survival is
spread over the merge stages. The first draft of this revision offered a
counterexample in which the first stage keeps 0.245 of its input and
every later stage keeps all of it. That counterexample is withdrawn:
every flushed byte enters stage 0, so it drops 0.755 bytes per byte
written, while only \(g_{\text{flow}} = 0.278\) bytes of garbage exist
per byte written. It met the floor because its drops are 0.278 of its
\emph{total} compaction input (2.71 per byte written); the floor sees
only total input and cannot tell the two apart, which is why it bounds
nothing about where survival goes. The bound that holds in the stage
model is the following.

\textbf{Proposition B.1″ (the drop budget bounds the elision saving;
stage model).} Fix a window in which A7 holds. Model the tree as merge
stages \(0..L-1\) with the accounting of Lemma D.7 and no trivial moves:
bytes that are not dropped pass through the stages in order; a
stage-\(i\) merge writes \(w_i = 1 + o_i\) bytes per source byte, less
one byte for each byte it drops, from the source or from the overlap;
the \(o_i\) are held fixed, as Theorem A.2(ii) holds its weights. Let
\(V\) be the bytes entering stage 0 in the window (the flushed bytes),
and \(G\) the obsolete bytes available to compaction in it: garbage
resident at the window\textquotesingle s start plus obsolete versions
flushed during it. Against the same flow with nothing dropped, the
compaction write cost saved is at most \(G\,(1 + \sum_{i=1}^{L-1}w_i)\),
so the relative saving is at most
\[\frac{G}{V}\cdot\frac{1 + \sum_{i\ge1}w_i}{\sum_{i\ge0}w_i} \;\le\; \frac{G}{V},\]
the last step because \(w_0 \ge 1\). With equal weights \(w\) the factor
is \((1 + (L-1)w)/(Lw)\), between \((L-1)/L\) and 1.

\emph{Proof.} Let \(D_j\) be the bytes dropped by stage-\(j\) merges. In
steady state a level releases what reaches it, so the volume entering
stage \(i\) is \(V_i = V - \sum_{j<i}D_j\), and by Lemma D.7 stage \(i\)
writes \(w_iV_i - D_i\). With nothing dropped the cost is
\(V\sum_iw_i\), so the saving is
\(\sum_jD_j\,(1 + \sum_{i>j}w_i) \le \big(\sum_jD_j\big)(1 + \sum_{i\ge1}w_i)\).
Every obsolete version is dropped at most once, so \(\sum_jD_j \le G\).
\(\blacksquare\)

\emph{Scope.} This is a statement about the stage model. Three things
lie outside it. The comparison flow keeps every stale version at the
bottom, so \(B_L\) and \(o_{L-1}\) would grow, and the model holds them
fixed. Trivial moves write nothing: a heavy share at stage 0 would make
the effective \(w_0\) fall below 1, and a heavy share at later stages,
as over a window that contains the load (Lemma D.7), lowers the no-drop
cost, so every relative figure below assumes that trivial moves are
rare; with up to 71\% of the bytes leaving levels \(\ge 1\) moving
trivially, as on the uniform whole run, a bound valid for any placement
of drops can exceed 58\%. A window that contains the load is not steady.

\emph{Values on \texttt{Assoc}.}

\begin{itemize}
\tightlist
\item
  \textbf{Whole run, load included} (the view in which
  \(S_{\text{flow}}\) is measured): memtable drops remove the same bytes
  from \(G\) and \(V\), so \(G/V \le g_{\text{flow}} = 0.278\). The
  window is not steady: the load grows the tree from empty, and bytes
  that end in L1--L7 never cross the later stages. Because levels
  \(1..L-1\) end at target in both flows, the saving is still at most
  \(G(1 + \sum_{i\ge1}w_i)\), but the no-drop cost is
  \(\sum_iw_i(V - R_{\le i})\), with \(R_{\le i}\) the bytes resident in
  levels \(1..i\) at the end. With that growth counted and no trivial
  moves, the bound is about 26--32\% of whole-run \(W - 1\) at T = 2
  (27--29\% at T = 6, 25--32\% at T = 10) for overlap constants from 1
  to 0, about half the 58\% of the 2026-09-11 formula at T = 2.
\item
  \textbf{Measured phase} (the view the objective scores, A8): every Put
  overwrites a live key, so each user byte creates an obsolete byte, and
  the tree the load leaves holds almost none, since its keys are unique.
  \(G/V\) is then close to 1 and so is the bound: it says little. Native
  compaction also spends most of the budget. Resident \(S\) of
  1.03--1.09 at run end leaves 0.09--0.27 GB of the
  phase\textquotesingle s 1.17 GB of obsolete bytes, so native drops
  77--92\% of them.
\end{itemize}

What binds in the measured phase is \emph{where} drops happen. A stale
version can be dropped only in a merge that also holds a newer version
of its key, and on \texttt{Assoc} the loaded versions sit in the deepest
levels (at T = 2, L6--L8 hold about 84\% of them and L8 alone about a
third). So a first overwrite drops its loaded version deep, while repeat
overwrites of hot keys can meet high: merge survival \(X/(S+O)\) at L1
was 0.84, 0.86 and 0.93 at T = 2, 6 and 10 (D-4). Against native rather
than against no drops, a controller can add at most the garbage native
leaves, \(G_{\text{res}} = (S-1)N_{\text{live}}\), and can otherwise
only move native\textquotesingle s drops higher. In the stage model its
saving over native is therefore at most
\[\sum_j D^{\text{nat}}_j\sum_{i=1}^{j}w_i \;+\; G_{\text{res}}\Big(1 + \sum_{i\ge1}w_i\Big),\]
with \(D^{\text{nat}}_j\) native\textquotesingle s dropped bytes per
level (§2 item 4). It follows from
\(\sum_jD^\pi_j \le G = \sum_jD^{\text{nat}}_j + G_{\text{res}}\), with
every drop placed at stage 0 where it saves the most. This ignores
locality, so it is an upper bound, and it is computable from the Hull-0
artifacts. A bound that uses locality needs a model of where overwrites
of one key meet, and is open. The proof also shows where elision pays
most: a byte dropped at stage \(j\) saves one byte there and every later
stage, so drops high in the tree are worth the most.

\textbf{Measured on \texttt{Assoc}} (D-3, measured denominator):
\(S_{\text{flow}} = 1.386\) at T = 2, 6 and 10;
\(g_{\text{flow}} = 0.278\); resident \(S\) 1.03--1.09. The
constant-survival \emph{estimate} --- not established by its proof, and
above what B.1″ allows for the whole run at every ratio when trivial
moves are rare --- is 58\%, 35\% and 35\% (audit §4). The B.1 and B.2
tables of 2026-09-11 (uniform workload, defective live-data estimate)
are withdrawn; Lemma D.14 replaces B.2\textquotesingle s space budget.

\textbf{Conjecture B.3′ (state-dependent value on a stationary
workload).} On a stationary workload there is a state-dependent trigger
policy \(\pi\) with
\(J_\beta(\pi) < \min_{\theta\in\Theta_s}J_\beta(\theta)\).

\emph{Status.} Unproven. The 2026-09-11 revision asserted it in its
commentary on Proposition B.3 (now Proposition H.3); B.3 does not imply
it. No learned or hand-written policy on record lies below the static
frontier (audit §2). The candidate mechanisms are items (b)--(d) of
Pathway A §4. They are tested without a learner at Gate N3; if none
works, the claim narrows to suite robustness and changing workloads
(Global acceptance).

\textbf{Definition (phase-adaptivity gap), retained.}
\[\mathcal G = \min_{\theta\in\Theta_s}J_\beta(\theta) - \sum_p\frac{n_p}{N}\min_{\theta\in\Theta_s}J_{\beta,p}(\theta),\]
the excess cost of the best single static configuration over a hindsight
oracle that switches static configuration per phase, where \(n_p/N\) is
phase \(p\)\textquotesingle s share of the measured operations; time
shares would depend on each configuration\textquotesingle s speed.

\textbf{Corollary B.4 (retained).} \(\mathcal G > 0\) is sufficient
evidence that adaptivity has value on that workload. \(\mathcal G = 0\)
is not evidence that it has none, because \(\mathcal G\) sees only the
phase component of adaptivity. Corollary D.12 now predicts
\(\mathcal G > 0\) in advance for the L0 trigger whenever
phases\textquotesingle{} best triggers differ.

\subsubsection{§4 Acceptance}\label{4-acceptance}

{\def\LTcaptype{none} % do not increment counter
\begin{longtable}[]{@{}
  >{\raggedright\arraybackslash}p{(\linewidth - 6\tabcolsep) * \real{0.0290}}
  >{\raggedright\arraybackslash}p{(\linewidth - 6\tabcolsep) * \real{0.3484}}
  >{\raggedright\arraybackslash}p{(\linewidth - 6\tabcolsep) * \real{0.4356}}
  >{\raggedright\arraybackslash}p{(\linewidth - 6\tabcolsep) * \real{0.1670}}@{}}
\toprule\noalign{}
\begin{minipage}[b]{\linewidth}\raggedright
\#
\end{minipage} & \begin{minipage}[b]{\linewidth}\raggedright
Criterion
\end{minipage} & \begin{minipage}[b]{\linewidth}\raggedright
Threshold
\end{minipage} & \begin{minipage}[b]{\linewidth}\raggedright
Instrument
\end{minipage} \\
\midrule\noalign{}
\endhead
\bottomrule\noalign{}
\endlastfoot
B-1 & Skew raises resident garbage (retained) & garbage fraction under
skew exceeds a skew-free control by ≥ 10 pp, on the measured denominator
& evaluator \\
B-3 & Survival is measurable and responds (retained) & \(\rho_i\) and
\(\eta_i\) per level, trivial moves excluded; lower under holding than
under native at the held level & compaction log \\
B-4 & Shift produces a phase-adaptivity gap (retained) &
\(\mathcal G > 0\) with the 95\% interval excluding 0 & hindsight-oracle
runner \\
B-5 & Deletes activate compensated size (retained) & non-zero
tombstone-driven file selections in the RocksDB LOG & LOG parse \\
WL-1 & Garbage is reported where it acts (replaces B-2) & per-level
\(\rho_i\) and dropped bytes per workload and mode & compaction log \\
WL-2 & The suite exercises every mode & for each mode, at least one
workload whose best static setting differs from the balanced-mode best &
Gate N2 \\
\end{longtable}
}

B-2 is retired: its threshold used Theorem B.1\textquotesingle s
ceiling, which its proof does not establish.

\begin{center}\rule{0.5\linewidth}{0.5pt}\end{center}

\subsection{Pathway C --- Comparator}\label{pathway-c--comparator}

\subsubsection{§0 Purpose}\label{0-purpose-2}

Compare the controller with the best static setting for the same priced
cost, measured on the same workload and binary, in the same session. The
claim is then about a class of static settings rather than one tuned
point.

\subsubsection{§1 The static class and the
comparator}\label{1-the-static-class-and-the-comparator}

\textbf{\(\Theta_s\) (fixed in advance, §0.6):}

\begin{itemize}
\tightlist
\item
  L0 trigger \(\in \{2, 4, 8\}\), restricted to admissible values
  (A-Impl-6), so that triggers the byte branch makes identical are not
  counted twice;
\item
  base size \(\in \{8, 16, 32\}\) MiB;
\item
  static multiplier profiles: uniform 1; the equal-fanout or
  survival-weighted profile at the measured \(c\) and
  \(v_i = a_i - t_i\) (Theorem A.2); a last-level-emptying profile that
  holds upper levels at 1.5--2 times from load (audit §3); uniform 0.75;
\item
  \texttt{compaction\_pri\ =\ kMinOverlappingRatio}, pinned;
\item
  \(T \in \{2, 6, 10\}\), with 14 and 20 at one point for the
  cross-\(T\) check.
\end{itemize}

\textbf{Comparator.} For each workload \(w\) and mode \(\beta\),
\(\theta^\star_\beta(w) = \arg\min_{\theta\in\Theta_s}J_\beta(\theta, w)\),
from measured, seed-paired runs.

\subsubsection{§2 Theory}\label{2-theory}

\textbf{Proposition C.1 (comparator validity; retained, restated for
priced cost).} If some \(\theta \in \Theta_s\) has priced cost at most
\(\pi\)\textquotesingle s on every term, then a static configuration is
at least as good as \(\pi\) in every mode, and \(\pi\)\textquotesingle s
contribution is nil whatever it achieves against any single setting.

\emph{Proof.} Immediate from dominance, since every \(J_\beta\) has
positive weights. \(\blacksquare\)

\textbf{Corollary C.2 (retained).} "\(\pi\) beats the native setting" is
strictly weaker than "\(\pi\) beats \(\theta^\star_\beta\)".

\textbf{Corollary C.3 (steady-state gains belong to the static class;
extended).} The steady-state optima of Proposition D.11 (\(K_0\) at
fixed prices), Proposition D.13 (interior profile, depth) and Theorem
A.2 (survival-weighted profile) are static settings. Any gain they
explain belongs to \(\Theta_s\), not to the controller.

\textbf{Proposition C.4 (one hull serves every mode).} For every
\(\beta > 0\), \(\theta^\star_\beta\) is a vertex of the lower convex
hull of the static points\textquotesingle{} priced-cost vectors
\((\mathcal C_W, \mathcal C_R, \mathcal C_S)\). The hull extracted once
per workload therefore contains every mode\textquotesingle s comparator.

\emph{Proof.} Proposition D.5 applied to the finite set \(\Theta_s\).
\(\blacksquare\)

\textbf{Definition C.5 (regret and suite robustness).} The regret of
\(\pi\) on workload \(w\) in mode \(\beta\) is
\[\text{Reg}_\beta(\pi, w) = \frac{J_\beta(\pi, w)}{J_\beta(\theta^\star_\beta(w), w)} - 1.\]
\(\pi\) is \emph{suite-robust} in mode \(\beta\) if
\[\max_w \text{Reg}_\beta(\pi, w) < \min_{\theta\in\Theta_s}\max_w\text{Reg}_\beta(\theta, w),\]
i.e. its worst regret across the suite is below the best worst-case
regret of any single static setting. This is where a min--max criterion
properly belongs. It is the regret form of robust tuning under workload
uncertainty {[}Endure{]}.

A controller can be suite-robust without beating
\(\theta^\star_\beta(w)\) on any single workload: it only has to avoid
being far from the best static setting everywhere, which no one static
setting may manage.

\subsubsection{§3 Drift and repeats}\label{3-drift-and-repeats}

\begin{itemize}
\tightlist
\item
  \textbf{Same session.} Comparator arms run in the same session as
  policy arms, in interleaved order, and the session id is recorded. The
  same static setting wrote 2.9--4.1\% more at T=2 in later sessions
  (audit).
\item
  \textbf{Repeats.} For the paired difference \(d\) between an arm and
  its comparator, the interval half-width is
  \(t_{n-1}\sigma_d/\sqrt n\), so the repeat count per cell is the
  smallest \(n\) with
  \[n \ge \left(\frac{t_{0.975,\,n-1}\,\sigma_d}{h}\right)^2,\] where
  \(\sigma_d\) is estimated at Gate N2 and \(h\) is the half-width
  required --- a quarter of the gap between \(\theta^\star_\beta\) and
  the next-best static point, or an effect size fixed in advance.
\end{itemize}

\subsubsection{§4 Acceptance}\label{4-acceptance-1}

{\def\LTcaptype{none} % do not increment counter
\begin{longtable}[]{@{}
  >{\raggedright\arraybackslash}p{(\linewidth - 6\tabcolsep) * \real{0.0402}}
  >{\raggedright\arraybackslash}p{(\linewidth - 6\tabcolsep) * \real{0.3293}}
  >{\raggedright\arraybackslash}p{(\linewidth - 6\tabcolsep) * \real{0.4820}}
  >{\raggedright\arraybackslash}p{(\linewidth - 6\tabcolsep) * \real{0.1285}}@{}}
\toprule\noalign{}
\begin{minipage}[b]{\linewidth}\raggedright
\#
\end{minipage} & \begin{minipage}[b]{\linewidth}\raggedright
Criterion
\end{minipage} & \begin{minipage}[b]{\linewidth}\raggedright
Threshold
\end{minipage} & \begin{minipage}[b]{\linewidth}\raggedright
Instrument
\end{minipage} \\
\midrule\noalign{}
\endhead
\bottomrule\noalign{}
\endlastfoot
C-1 & Static class sampled (retained, extended) & every knob level of
\(\Theta_s\) present; ≥ 4 hull vertices per (workload, \(T\)) & sweep
output \\
C-2 & Hull points decidable (retained) & repeats until the interval
width is under half the spacing between neighbouring hull points &
paired evaluator \\
CMP-3 & Per-workload comparison &
\(J_\beta(\pi) - J_\beta(\theta^\star_\beta)\) with paired interval, per
(workload, mode, \(T\)) & paired evaluator \\
C-6 & Cross-\(T\) (retained) & comparison repeated with
\(\theta^\star_\beta\) taken over \(T \in \{2, 6, 10, 14, 20\}\) &
cross-\(T\) sweep \\
CMP-7 & Suite robustness & Definition C.5 reported per mode & suite
evaluator \\
CMP-8 & Same-session twins & every comparison paired within one session,
session id recorded & run manifest \\
\end{longtable}
}

The 2026-09-11 criteria C-3, C-4 and C-5 are historical (history §18.1).

\begin{center}\rule{0.5\linewidth}{0.5pt}\end{center}

\subsection{Pathway E --- Guard (Programme 2; on
hold)}\label{pathway-e--guard-programme-2-on-hold}

Programme 1 has no service-level objective, so it has no guard.
RocksDB\textquotesingle s own slowdown and stop triggers and
pending-bytes limits are the only safety mechanism, together with the
multiplier and trigger bounds of A-Impl-6 and A-Impl-7.

\textbf{Retained for Programme 2:} the 2026-09-11 Proposition E.1 (the
realised policy is a mixture of the policy and its projection onto the
shield\textquotesingle s action set, so a shield that fires often
replaces the policy) and Corollary E.2 (the shield\textquotesingle s
influence is bounded by its override probability, and a shield action
set narrower than the policy\textquotesingle s biases every override the
same way), the criteria E-1 to E-5 as defined there, and three
calibration lessons:

\begin{itemize}
\tightlist
\item
  limits must be calibrated in the unit the criterion scores --- frames,
  not episodes (the inspection-paradox finding of 2026-09-14);
\item
  every term of the force condition must be calibrated, not only the
  per-level ones;
\item
  calibration windows must exclude the post-load transient. E-5 was
  6--17 times its limit over whole runs and 0--1.4\% after the first 30
  s (audit).
\end{itemize}

Verdicts to date are in history §18.1.

\begin{center}\rule{0.5\linewidth}{0.5pt}\end{center}

\subsection{Pathway F --- Phase-aware objective switching (Programme
2)}\label{pathway-f--phase-aware-objective-switching-programme-2}

\textbf{Status:} second programme. Its 2026-09-11 description --- a
layer above the controller that switches a parameter set per detected or
forecast phase; the F-detect and F-forecast versions; the leak guard;
switchable knobs only --- is retained. Three changes:

\begin{itemize}
\tightlist
\item
  Under Programme 1, what switches per phase is the priority vector
  \(\beta\) (and \(K_0\) at its phase optimum, Corollary D.12), not an
  SLO manifest. SLO manifests return with Programme 2.
\item
  Corollary D.12 predicts \(\mathcal G > 0\) for the L0 trigger in
  advance whenever the phases\textquotesingle{} best triggers differ.
\item
  The audit\textquotesingle s suggestion of a slow tuner above native
  RocksDB is this layer.
\end{itemize}

\textbf{Definition (three controllers over one schedule; retained).} For
phases \(p = 1..P\) with boundaries known to the evaluator:
\(J_{\text{oracle}}\) switches the best static setting per phase exactly
at the boundaries (its advantage over the best single static setting is
\(\mathcal G\)); \(J_{\text{react}}\) chooses from the current
window\textquotesingle s measured mix, with no look-ahead;
\(J_{\text{pred}}\) chooses from a predicted class for the next episode,
with lead time \(\ell\).

\textbf{Proposition F.1 (decomposition; retained).}
\(\mathcal G = (J_{\text{static}} - J_{\text{react}}) + (J_{\text{react}} -
J_{\text{pred}}) + (J_{\text{pred}} - J_{\text{oracle}})\): the value of
reacting, the value of anticipating, and the remaining loss.

\emph{Proof.} Telescoping. \(\blacksquare\)

\textbf{Proposition F.2 (catch-up bounds the value of anticipation;
corrected).} Suppose \(J_{\text{react}}\) and \(J_{\text{pred}}\) choose
the same settings except that the predictor switches \(\ell\) earlier at
each transition. Let \(e_r(t) \ge 0\) be the reactive
controller\textquotesingle s excess cost rate over the new
setting\textquotesingle s steady cost during the first
\(\Delta_{\text{shape}}\) after a transition (the catch-up), and assume
the new setting\textquotesingle s steady cost is a floor for both
controllers after the transition, and that running the new setting early
costs the predictor \(w(t) \ge 0\) before the transition. Then, per
transition,
\[J_{\text{react}} - J_{\text{pred}} \;\le \int_0^{\Delta_{\text{shape}}} e_r(t)\,dt.\]
The difference can be negative if the pre-transition cost outweighs the
saving. A predictor reaches the bound only if it completes the reshaping
before the transition (\(\ell \ge \Delta_{\text{shape}}\)) at no
pre-transition cost.

\emph{Proof.} The two controllers differ only in
\([-\ell, \Delta_{\text{shape}}]\) around a transition. There
\(J_{\text{react}} - J_{\text{pred}} = \int e_r -
\int e_p - \int w\), with \(e_p \ge 0\) by the floor assumption and
\(w \ge 0\), so the difference is at most \(\int e_r\). \(\blacksquare\)

\emph{Correction.} The 2026-09-11 statement bounded the anticipation
term by the cost over the first
\(\max(0, \Delta_{\text{shape}} - \ell)\) of the read phase. That
quantity is what a predictor with lead \(\ell\) still leaves --- part of
the third term of F.1 --- not the anticipation term. The conclusion
stands: on a gradual transition the catch-up excess is small, so the
anticipation term is small; transition abruptness is the axis to sweep.

\textbf{Acceptance (retained, Programme 2):} F-1 (\(\mathcal G > 0\);
this is B-4), F-2 (reacting captures a measured share of
\(\mathcal G\)), F-3 (anticipation adds value on the abrupt schedule),
F-4 (whole-run bounds hold under switching), F-5 (no churn), F-6
(forecaster is honest), F-7 (lead-time dependence shown), as defined in
the 2026-09-11 revision. F-4\textquotesingle s bounds are re-stated when
Programme 2\textquotesingle s constraints are.

\begin{center}\rule{0.5\linewidth}{0.5pt}\end{center}

\subsection{Execution order (Programme
1)}\label{execution-order-programme-1}

Each gate\textquotesingle s cost is stated in runs. Node-hours follow
once Gate N1 fixes run length.

\textbf{Gate N0 --- code, instruments and tests (node time only for
builds, tier-2 tests and the preflight).}

\begin{enumerate}
\def\labelenumi{\arabic{enumi}.}
\tightlist
\item
  \textbf{(Done 2026-09-25, \texttt{1fc8852}.)} Commit the D-10 to D-12
  work, the D-12 verdict and the audit (audit §5, item 1).
\item
  \texttt{level\_target\_multipliers} with its tests (ACT-1; the
  implementation plan is
  \texttt{docs/IMPLEMENTATION\_PLAN\_PROGRAMME1.md}).
\item
  Per-level read counters, keyed by the version\textquotesingle s level
  and aggregated across threads, with each level\textquotesingle s block
  reads split into the hit\textquotesingle s read and false-positive
  reads (D §4), and \(R_{blk}\); the slot-blocking spans (L0 due but
  waiting, and the start level of the job holding the slot) with the
  operations served in them, for the slot-blocking charge (D §4); the
  active memtable\textquotesingle s fill at every L0 decision (H §2);
  \(\rho_i\), \(o_i\), dropped bytes and trivially moved bytes \(t_i\)
  per source level, per compaction; turnover logging; per-release slot
  wait (from the later of a level\textquotesingle s score reaching 1 and
  its previous job ending, to the job starting) and slot occupancy, for
  \(\hat\omega_j\), \(\omega_j\) and \(\zeta\) (G.4); at every decision,
  the score order, the running compaction\textquotesingle s start level
  and the pending-compaction estimate with the multipliers applied
  (A-Impl-3), read in-process from the current version, for the
  queue-position and backlog inputs (G §3). Repair the existing
  per-level telemetry counters \texttt{trivial\_move\_jobs} and
  \texttt{trivial\_move\_bytes}, which read zero:
  \texttt{NotifyOnCompactionCompleted} passes
  \texttt{Compaction::is\_trivial\_move()}, a flag only the universal
  picker sets, and the trivial-move branch of
  \texttt{BackgroundCompaction} never fills
  \texttt{total\_input\_bytes}. \(\hat\xi_j\) (G §3) must not be read
  from them until then.
\item
  The settle step after the load (\texttt{WaitForCompact}, then the
  \(h_w\) hold) and the measured-phase stamp at the first
  \texttt{mixgraph} operation (A8, H §5).
\item
  Evaluator: \(\mathcal C_W\), \(\mathcal C_R\), \(\mathcal C_S\) and
  \(J_\beta\) over the measured phase, from operation \(n_w\) to the end
  of the drain --- a ticker snapshot and an internal-stats stall
  snapshot at \(n_w\), \(R_{blk}\) from
  \texttt{rocksdb.bloom.filter.full.positive}, compactions windowed by
  completion time, and \(H\) sampled with the operation count at every
  version install on every arm, native included. Self-check on every
  arm: the per-interval operation counts sum to the measured operations,
  and the last sampled \(H\) equals \texttt{rocksdb.live-sst-files-size}
  at the end of the drain. And the OBJ-6 diagnostics.
\item
  The C++ inference plugin, weight push, masked target and fallback (H
  §4, H §6).
\item
  Device price calibration: \(c_w\), \(c_f\), \(c_{blk}\), \(c_{sk}\),
  \(c_s\) (OBJ-2).
\item
  The dated \texttt{PREREGISTRATION.md} entries of §0.6.
\item
  Test suites for every item above, and the preflight (CLAUDE.md
  "Tests"; plan §6). Every later gate that needs a long node run starts
  only with a preflight marker matching the current \texttt{db\_bench},
  plugin and code hashes.
\end{enumerate}

\textbf{Gate N1 --- admission test on pilot native runs (node time:
about 18 runs).} Amended 2026-10-01 (PREREGISTRATION D-16): the
2026-09-11 programme\textquotesingle s artifacts are out of scope
(owner, 2026-09-30) and carry no host log, so Gate N1 runs its own pilot
\texttt{native} arms at the default point, three per (workload, \(T\)).
PROP-1 for the levels the pilots decide. It fixes the run length: the
measured phase must contain at least \(n_{\text{turn}}\) turnovers of
the deepest pooled level, with \(n_{\text{turn}}\) = 10 (D-16). A cell
with no pool applies the same rule to its deepest interior level with
enough turnovers, L2. The load stays fixed and only \texttt{mixgraph}
grows, so the tree the pools were decided on is the tree every later run
settles. Amended 2026-10-02 (PREREGISTRATION D-19): a candidate the
pilots leave undecided is not pooled, does not set the run length, and
is not decided later on Gate N2. \textbf{Status (D-19):} run on
2026-10-01. No cell has a pool, so the run length comes from L2 (D-19 §3
estimates the rungs).

\textbf{Gate N2 --- static comparator on the new binary.} ACT-4 (native
parity) first. Then \(\Theta_s\) on \texttt{Assoc} and one Zipfian
workload at T = 2, 6 and 10, in the same-session design, with repeat
counts from C §3. Runs: \(|\Theta_s| \times\) workloads \(\times\) \(T\)
values \(\times\) repeats.

\textbf{Gate N3 --- learner-free test of Conjecture B.3′.} Hand-written
state rules --- \(K_0\) tracking Proposition D.11 at measured prices;
releasing on neighbour fill; holding for garbage; L0 compacting early
while the slot is idle and reads are heavy (Pathway A §4 (e)); interior
levels holding releases while L0 is due or within one flush of due
(Pathway A §4 (f)) --- against \(\theta^\star_\beta\) in each mode, with
the stall rule applied (Global acceptance). This decides whether the
stationary-workload claim stays before any learner run. If neither slot
rule beats \(\theta^\star_\beta\) in read priority, read
priority\textquotesingle s claim narrows to changing workloads and suite
robustness (claims 2 and 3) before Gate N5.

\textbf{Gate N4 --- learner smoke test.} One workload, read priority,
one \(T\): ARCH-1 to ARCH-6, OBJ-1, OBJ-3, OBJ-4.

\textbf{Gate N5 --- mode matrices.} Read priority across the suite; then
write and space priority (only \(\beta\) changes), with the same
instruments. The claim is recorded before N5 (Global acceptance).

\textbf{Gate N6 --- ablations.} Propagation on and off (PROP-4);
counterfactual charge on and off; neighbour charge from the full value
versus the attributed-cost head (§0.6 item 11); prior on and off; plugin
versus socket inference (the plugin\textquotesingle s remote-inference
mode, same policy with Q evaluated in the trainer).

Programme 2 (Pathways E and F) starts after Programme
1\textquotesingle s claim is recorded.

\begin{center}\rule{0.5\linewidth}{0.5pt}\end{center}

\subsection{Global acceptance}\label{global-acceptance}

\textbf{Claims}, one or more of which is recorded in
\texttt{PREREGISTRATION.md} before Gate N5:

\begin{enumerate}
\def\labelenumi{\arabic{enumi}.}
\tightlist
\item
  \textbf{Per-workload gain.} In mode \(\beta\),
  \(J_\beta(\pi) < J_\beta(\theta^\star_\beta)\) with the paired
  interval below zero on named workloads (CMP-3), and against the
  cross-\(T\) comparator (C-6).
\item
  \textbf{Suite robustness.} Definition C.5 holds in the mode (CMP-7).
\item
  \textbf{Changing workloads.} \(\mathcal G > 0\) (B-4) and the
  controller captures a measured share of it (F-2 form).
\item
  \textbf{Propagation contributes.} PROP-2 to PROP-4, claimed only for
  (workload, \(T\)) cells whose pool holds at least two admitted levels
  (G §4, scope decision): T=2 on \texttt{Assoc} at the current size, T=6
  if the admission test admits two levels, never T=10 at this size. Gate
  N1 (D-19) found no pool in any cell, so this claim is not made at the
  current size.
\end{enumerate}

\textbf{Stall rule} (applies to every claim). \(J_\beta\) prices no time
(D §1), so a policy could lower it by deferring work until writes stall.
A claim therefore also requires, on the same paired runs and reported
beside \(J_\beta\) rather than priced into it (OBJ-6), measured-phase
stall seconds no more than \(\delta_{\text{stall}}\) above the
comparator\textquotesingle s and foreground throughput no more than
\(\delta_{\text{thr}}\) below it. Both use the paired
interval\textquotesingle s bound; both margins are fixed in advance
(§0.6 item 12). Stall seconds are read as in D-11, from the
internal-stats \texttt{Cumulative\ stall} line differenced over the
measured phase, never from \texttt{rocksdb.stall.micros}.

\textbf{If none holds.} If Gate N3 finds no state-dependent rule that
beats \(\theta^\star_\beta\) and Gate N5 confirms it, the paper is an
analysis paper: the structure of the static optimum (Propositions D.11
and D.13, Theorem A.2), depth monotonicity (Lemma A.6), the corrected
elision analysis (Propositions B.1′ and B.1″), the cost of exploration
(Proposition H.3), and the propagation measurements (Proposition G.1).

\begin{center}\rule{0.5\linewidth}{0.5pt}\end{center}

\subsection{References}\label{references}

\begin{itemize}
\tightlist
\item
  {[}Cao et al.{]} Z. Cao, S. Dong, S. Vemuri, D. H. C. Du.
  Characterizing, Modeling, and Benchmarking RocksDB Key-Value Workloads
  at Facebook. USENIX FAST 2020.
\item
  {[}Cooper et al.{]} B. F. Cooper, A. Silberstein, E. Tam, R.
  Ramakrishnan, R. Sears. Benchmarking Cloud Serving Systems with YCSB.
  ACM SoCC 2010.
\item
  {[}Cosine{]} S. Chatterjee, M. Jagadeesan, W. Qin, S. Idreos. Cosine:
  A Cloud-Cost Optimized Self-Designing Key-Value Storage Engine. PVLDB
  15(1), 2021.
\item
  {[}Das \& Dennis{]} I. Das, J. E. Dennis. A Closer Look at Drawbacks
  of Minimizing Weighted Sums of Objectives for Pareto Set Generation in
  Multicriteria Optimization Problems. Structural Optimization
  14:63--69, 1997.
\item
  {[}Dostoevsky{]} N. Dayan, S. Idreos. Dostoevsky: Better Space-Time
  Trade-Offs for LSM-Tree Based Key-Value Stores via Adaptive Removal of
  Superfluous Merging. ACM SIGMOD 2018.
\item
  {[}Endure{]} A. Huynh, H. A. Chaudhari, E. Terzi, M. Athanassoulis.
  Endure: A Robust Tuning Paradigm for LSM Trees Under Workload
  Uncertainty. PVLDB 15(8), 2022.
\item
  {[}Foerster et al.{]} J. N. Foerster, G. Farquhar, T. Afouras, N.
  Nardelli, S. Whiteson. Counterfactual Multi-Agent Policy Gradients.
  AAAI 2018.
\item
  {[}How to Grow{]} How to Grow an LSM-tree? Towards Bridging the Gap
  Between Theory and Practice. ACM SIGMOD 2025. Source of the M3
  convention used in P0-7.
\item
  {[}Kearns \& Singh{]} M. Kearns, S. Singh. Near-Optimal Reinforcement
  Learning in Polynomial Time. Machine Learning 49:209--232, 2002.
  Simulation lemma.
\item
  {[}Miettinen{]} K. Miettinen. Nonlinear Multiobjective Optimization.
  Kluwer, 1999.
\item
  {[}Mnih et al.{]} V. Mnih et al. Human-Level Control through Deep
  Reinforcement Learning. Nature 518:529--533, 2015.
\item
  {[}Monkey{]} N. Dayan, M. Athanassoulis, S. Idreos. Monkey: Optimal
  Navigable Key-Value Store. ACM SIGMOD 2017.
\item
  {[}Ng et al.{]} A. Y. Ng, D. Harada, S. Russell. Policy Invariance
  under Reward Transformations: Theory and Application to Reward
  Shaping. ICML 1999.
\item
  {[}O\textquotesingle Neil et al.{]} P. O\textquotesingle Neil, E.
  Cheng, D. Gawlick, E. O\textquotesingle Neil. The Log-Structured
  Merge-Tree (LSM-Tree). Acta Informatica 33(4):351--385, 1996.
\item
  {[}RusKey{]} D. Mo, F. Chen, S. Luo, C. Shan. Learning to Optimize
  LSM-trees: Towards a Reinforcement Learning Based Key-Value Store for
  Dynamic Workloads. ACM SIGMOD 2024; arXiv:2308.07013. Lerp and policy
  propagation in §5; evaluation in §7.
\item
  {[}Sarkar et al.{]} S. Sarkar, D. Staratzis, Z. Zhu, M. Athanassoulis.
  Constructing and Analyzing the LSM Compaction Design Space. PVLDB
  14(11), 2021.
\item
  {[}van Hasselt et al.{]} H. van Hasselt, A. Guez, D. Silver. Deep
  Reinforcement Learning with Double Q-Learning. AAAI 2016.
\item
  {[}Wolpert \& Tumer{]} D. H. Wolpert, K. Tumer. Optimal Payoff
  Functions for Members of Collectives. Advances in Complex Systems
  4(2--3):265--279, 2001.
\item
  {[}Zhu et al.{]} Z. Zhu, J. H. Mun, A. Raman, M. Athanassoulis.
  Reducing Bloom Filter CPU Overhead in LSM-Trees on Modern Storage
  Devices. DaMoN 2021.
\end{itemize}

RocksDB behaviour (score computation, file picking, \texttt{SetOptions})
is cited to the pinned source, not to documentation.

\end{document}
